\documentclass[UTF8,12pt,a4paper,fontset=fandol]{ctexart}

% 支持从项目根目录、深度学习目录或本文件所在目录加载共享样式。
\makeatletter
\def\input@path{{../}{../../}}
\makeatother
\usepackage{researchnotes}

\title{生成扩散模型的原理与前置知识}
\author{}
\date{}

\begin{document}
\maketitle
\begin{abstract}
生成扩散模型把复杂分布的学习分解为不同噪声尺度下的统计推断，再通过随机过程或确定性动力系统完成采样。本文从概率分布、条件期望、高斯运算、信息论、变分推断与神经网络优化出发，完整推导离散扩散的前向边缘、后验分布、变分目标及噪声预测，建立去噪、得分估计和连续时间动力学之间的联系，进一步说明采样加速、条件引导、潜空间建模、逆问题、流匹配及一致性学习。全文以可手算的例子解释抽象对象，区分精确恒等式、建模近似、数值误差与经验选择；所有算法均以数学关系和自然语言说明，不包含程序代码或伪代码。
\end{abstract}
\tableofcontents

\section{研究对象、基本直觉与学习路线}
\label{sec:diff-intro}
\subsection{为什么从噪声出发可以生成结构}
考虑一组手写数字图像。每幅图像都可以排列成一个很长的向量，但任意抽取各像素通常只会得到杂乱的斑点，因为真实图像的像素具有很强的依赖关系。生成模型（generative model）要学习这种联合变化规律，使随机抽取的新向量也呈现合理笔画。这里的任务不是为已有图片查找最相近的训练样本，也不是给每个像素独立选择最常见的颜色，而是近似一个高维概率分布。生成的随机性可以来自起点、路径中的扰动或最终观测步骤；只要整体过程规定了输出的概率规律，就构成一个生成模型。

\keyterm{扩散模型（diffusion model）}选择先研究一个容易描述的过程：把真实样本逐渐扰乱，使不同样本的分布越来越难以区分。加噪到足够强时，分布可接近预先指定的简单噪声分布。随后学习在每一种噪声强度下，怎样从当前状态向较少噪声的状态转移。生成时便可以抽取简单噪声，再逐步恢复数据的统计结构。前向扰动与反向生成是两种不同用途的过程，时间方向相反；“恢复”指恢复分布中的结构，不要求找回最初被破坏的某一幅具体图像。

设一维数据只可能取 $-2$ 和 $2$，两种情况各占一半。观察 $Y=X+\varepsilon$，其中 $\varepsilon\sim\mathcal N(0,1)$ 与 $X$ 独立。如果看到 $Y=2.5$，来自正端点的可能性更大；如果看到 $Y=0$，两种解释同样合理。此时最佳平方误差估计为零，但零本身并不是原始数据的可能值。这一现象已经包含扩散学习的重要区分：一次去噪的条件均值不一定是一个真实样本，而由许多小步骤构成的概率生成过程可以重建两个分离的峰。把去噪网络输出直接看成“已经生成的图片”，会忽略其所在噪声尺度和后续采样机制。

所谓“知道加入了哪一份噪声”，只是在构造训练样本时成立。研究者先抽取干净数据，再亲自抽取噪声，因此二者都可用作监督；网络实际接收的却只有加噪结果、时间以及可能的外部条件。不同的干净数据与噪声可以产生相同观测，网络通常不能唯一反演每次抽样的噪声。它学到的是关于未知量的统计估计。这解释了为什么训练标签可以精确获得，而总体最优回归误差仍不必为零，也解释了监督去噪为何能服务于无标签的生成建模。

\subsection{模型、目标、网络与采样器的分工}
理解扩散需要分开四个层次。模型规定联合分布或连续时间动力学；训练目标决定怎样从数据识别其中的未知函数；神经网络是表示这些函数的工具；采样器把已经学到的函数变成有限次计算的输出。同一个噪声预测网络可以配合不同采样器，同一个采样器也可以使用不同网络结构。改变步数通常改变数值近似，而改变噪声过程或预测参数化可能同时改变统计目标。因此，“某个模型只需要若干步”必须附带训练方式、采样器、误差容忍度和任务背景，不能单靠模型名称判断。

早期扩散概率建模、去噪得分学习、去噪扩散概率模型（denoising diffusion probabilistic model, DDPM）及随机微分方程表述提供了几条相互连接的研究路线\cite{sohl2015,song2019,ho2020,song2021sde}。本文围绕这些稳定的基础机制展开，不把某一时刻的排行榜或特定产品配置当成数学原理。基础恒等式会在文中推导；反向随机过程等需要随机分析的结论将明确给出条件和引用边界；架构、优化及评价中的选择则解释其动机，而不声称存在适用于一切数据的最优方案。

\subsection{前置知识与阅读顺序}
全文先建立概率和线性代数工具，随后从有限步马尔可夫链推导 DDPM，再把均方回归转写为得分学习。连续时间部分补充布朗运动、伊藤公式、福克尔--普朗克方程和连续性方程，使反向随机微分方程与概率流常微分方程有共同的解释。后半部分讨论数值采样、网络、训练诊断、条件引导、潜空间、逆问题及相邻方法，最后用综合例题检查知识链条。首次阅读可以先完成离散时间部分，再回到连续时间；但理解采样加速与流匹配时，应补上相应微分方程知识。

全文以 $X_0\in\mathbb R^d$ 表示干净随机向量，以 $x_0$ 表示其取值；$d$ 是数据维数，$I_d$ 为单位矩阵，$\|x\|^2=x^\mathsf T x$。训练数据为 $\mathcal D=\{x_0^{(i)}\}_{i=1}^{N}$，真实分布记作 $P_{\mathrm{data}}$。若存在密度，写成 $p_{\mathrm{data}}$；不要求原始分布一定有全维密度。离散时间为 $t=0,1,\ldots,T$，连续时间另用 $\tau\in[0,\mathcal T]$。$q$ 表示已指定的数据扰动联合规律及其条件分布，$p_\theta$ 表示含参数 $\theta$ 的生成模型；在连续部分，$p_\tau$ 专门表示真实扰动过程在时刻 $\tau$ 的密度。所有对数均为自然对数，除非明确换算为比特。

\section{概率论基础：观察、推断与条件期望}
\label{sec:diff-probability}
\subsection{分布、密度和高维随机性}
概率分布给事件分配概率，密度则是相对于参考测度的局部表示。若 $X$ 具有密度 $p$，则 $\Pr(X\in A)=\int_A p(x)\,\mathrm dx$。密度可以超过一，而单点概率仍为零；高维图像上某个样本密度的数值，也依赖像素使用何种单位。离散像素对应概率质量，连续归一化像素常用密度建模，二者不能在评价时直接混用。把整数像素加上区间内的小扰动称为去量化（dequantization），其作用是构造一个明确的连续观测模型，而不是凭空提高数据的信息量。

样本可以来自真实总体，也可以来自有限训练集的经验分布 $\widehat P_N=N^{-1}\sum_i\delta_{x_0^{(i)}}$，其中 $\delta_x$ 表示在 $x$ 处的点质量。经验分布始终可定义，即使所有样本位于低维曲面上。直接对这种分布取对数密度梯度通常没有意义，但与非退化高斯噪声卷积后会得到平滑、处处正的密度。这是加噪在数学上具有价值的原因之一：它不仅让图像变模糊，还为原本可能奇异的数据规律建立了可以微分的表示。

对于可积函数 $h$，期望是按概率权重求平均的运算 $\mathbb E[h(X)]$。若二阶矩有限，均值 $m=\mathbb E[X]$，协方差为 $\operatorname{Cov}(X)=\mathbb E[(X-m)(X-m)^\mathsf T]$。独立随机向量的协方差可相加，但相关向量还要包含交叉项。扩散推导反复使用独立性：新噪声独立于当前状态，累计噪声独立于原始数据。若噪声来自一个依赖输入的机制，这些简单公式不再原样成立，必须重新写出条件分布。

\subsection{贝叶斯公式与信息的方向}
联合密度可分解为 $p(x,y)=p(x)p(y\mid x)$。对 $p(y)>0$ 的观测，贝叶斯公式给出
\begin{equation}
 p(x\mid y)=\frac{p(y\mid x)p(x)}{p(y)},\qquad
 p(y)=\int p(y\mid x)p(x)\,\mathrm dx.
 \label{eq:diff-bayes}
\end{equation}
已知噪声机制意味着 $p(y\mid x)$ 已知，不意味着 $p(x\mid y)$ 同样已知；后验还依赖数据先验 $p(x)$。如果给同一幅模糊轮廓配上“手写数字”或“自然物体”这两种不同先验，对原图的合理解释就会变化。扩散模型学习的核心正是这种数据依赖性。已知加噪机制仍然需要训练，并不是算法故意把简单问题复杂化。

边缘化是对所有可能解释求和或积分，不是挑选一个最可能的解释。最大后验估计选择 $p(x\mid y)$ 最大的位置，条件均值计算 $\mathbb E[X\mid Y=y]$，后验抽样则随机选取一个符合后验权重的结果。在多峰情形下，三者可以完全不同。生成任务通常需要保留后验的不确定性，而平方回归最自然得到条件均值。扩散之所以能用回归来生成，需要靠整个反向转移或连续流把这些局部统计量组合起来，不能仅凭“去噪网络”四个字跳过这一层论证。

\subsection{平方误差为何得到条件均值}
\begin{proposition}[条件均值的最小平方误差性质]
\label{prop:diff-regression}
设随机向量 $U$ 满足 $\mathbb E\|U\|^2<\infty$，$V$ 为观测，$m(V)=\mathbb E[U\mid V]$。对任意使误差可积的可测函数 $h$，有
\begin{equation}
 \mathbb E\|U-h(V)\|^2
 =\mathbb E\|U-m(V)\|^2+\mathbb E\|m(V)-h(V)\|^2.
 \label{eq:diff-projection}
\end{equation}
因此总体均方误差的最优预测为条件均值，在 $V$ 的分布意义下几乎处处唯一。
\end{proposition}
\begin{proof}
把 $U-h(V)$ 分解为 $U-m(V)+m(V)-h(V)$ 并展开平方。交叉项的条件期望为 $\mathbb E[U-m(V)\mid V]^\mathsf T(m(V)-h(V))=0$。再利用全期望公式对 $V$ 求平均便得结论。这里的正交性是平均意义下的正交性，并不要求两个随机向量独立。
\end{proof}

命题\ref{prop:diff-regression}是后文噪声预测、干净数据预测、去噪得分匹配和流匹配的共同基础。它描述无限数据、足够大函数空间中的总体最优解，有限神经网络只能寻找可表示范围内的近似最优。训练误差与理论最小误差之间，还可能包含有限样本误差、优化误差和分布偏移。若给不同时间赋予严格正的权重，在函数能独立表示每个时间的理想情形下，逐时间的条件均值不变；但共享有限网络时，权重会改变不同任务之间的资源分配。

\begin{example}[双点数据的后验均值]
\label{ex:diff-two-point}
令 $X$ 等概率取 $-a,a$，$a>0$，并令 $Y=X+\sigma\varepsilon$，$\sigma>0$、$\varepsilon\sim\mathcal N(0,1)$。两个后验概率之比为 $\exp(2ay/\sigma^2)$，于是
\[
 \mathbb E[X\mid Y=y]=a\tanh\!\left(\frac{ay}{\sigma^2}\right).
\]
当 $y=0$ 时估计为零；当 $y$ 远大于噪声尺度时估计接近 $a$。较大噪声使两个解释的权重更接近，均值更容易落在两个峰之间。这是后验不确定性的表现，不等于网络没有学会原始数据的支持集。
\end{example}

全方差公式进一步说明去噪的极限：$\operatorname{Cov}(U)=\mathbb E[\operatorname{Cov}(U\mid V)]+\operatorname{Cov}(\mathbb E[U\mid V])$。前一项是观测后仍剩下的不确定性，后一项是不同观测之间可解释的变化。把条件均值当成随机样本会丢掉第一项，因而容易低估变化幅度。后面的后验方差推导会再次出现这一结构；它能帮助识别“预测均值正确，所以转移分布也正确”这一常见误解。

\section{线性代数、高斯分布与微积分工具}
\label{sec:diff-gaussian}
\subsection{向量几何与高斯密度}
图像可以保留高度、宽度和通道三个轴，也可以展平为 $d$ 维列向量。展平仅改变索引，不会自动删除空间关系；是否利用这些关系由网络结构决定。向量内积 $u^\mathsf T v$ 衡量方向关系，平方范数 $\|u\|^2$ 累计各坐标误差。若各通道的物理单位或数值尺度不同，未经处理的欧氏距离会偏向大尺度通道。因此，扩散中的标准化既影响神经网络优化，也改变“各方向同样强的噪声”所对应的几何含义。

对称正定矩阵 $\Sigma$ 满足任意非零 $u$ 都有 $u^\mathsf T\Sigma u>0$，因而具有正特征值、可逆且行列式为正。多元高斯分布（multivariate Gaussian distribution）的密度为
\begin{equation}
 \mathcal N(x;m,\Sigma)=\frac{\exp[-\tfrac12(x-m)^\mathsf T\Sigma^{-1}(x-m)]}
 {(2\pi)^{d/2}(\det\Sigma)^{1/2}}.
 \label{eq:diff-normal}
\end{equation}
协方差的特征向量决定等密度椭球的轴，特征值决定各方向的尺度。$\Sigma=\sigma^2I_d$ 时称为各向同性（isotropic），每个方向的噪声方差相同。若协方差奇异，分布仍可能存在，但不具有式\eqref{eq:diff-normal}所写的全维密度；不能把零行列式代入该式继续计算。

若 $Z\sim\mathcal N(0,I_d)$，则 $m+AZ$ 的均值为 $m$、协方差为 $AA^\mathsf T$。若 $Z_1,Z_2$ 独立且都为标准高斯，则 $aZ_1+bZ_2$ 与 $\sqrt{a^2+b^2}Z$ 同分布。这里强调的是分布相同，不是两边在同一次实验中逐点相等。推导直接加噪公式时，常把许多独立噪声合并成一个新的标准高斯；这个新变量不是某一步原来抽取的噪声，也不能在不同时间任意共享后仍声称得到原马尔可夫链的联合分布。

\subsection{联合高斯的条件分布与配方}
设 $(U,V)$ 联合高斯，均值分别为 $m_U,m_V$，协方差分块为 $\Sigma_{UU},\Sigma_{UV},\Sigma_{VU},\Sigma_{VV}$，且 $\Sigma_{VV}$ 正定。条件均值和协方差为
\begin{align}
 \mathbb E[U\mid V=v]&=m_U+\Sigma_{UV}\Sigma_{VV}^{-1}(v-m_V),\label{eq:diff-gauss-cond-mean}\\
 \operatorname{Cov}(U\mid V=v)&=\Sigma_{UU}-\Sigma_{UV}\Sigma_{VV}^{-1}\Sigma_{VU}.
 \label{eq:diff-gauss-cond-var}
\end{align}
可把 $U$ 分成对 $V$ 的线性预测与残差，检验残差和 $V$ 的协方差为零；联合高斯性进一步保证二者独立，从而得到上述条件分布。没有联合高斯假设时，线性预测仍有意义，但零相关不保证独立，条件均值也不一定是线性的。

另一种实用推导方法是配方。若未知向量 $u$ 的未归一化密度正比于 $\exp[-\tfrac12 u^\mathsf T A u+b^\mathsf T u]$，其中 $A$ 正定，则
\[
 -\tfrac12 u^\mathsf T A u+b^\mathsf T u
 =-\tfrac12(u-A^{-1}b)^\mathsf T A(u-A^{-1}b)+\tfrac12b^\mathsf T A^{-1}b.
\]
最后一项与 $u$ 无关，归入归一化常数，因此均值为 $A^{-1}b$，协方差为 $A^{-1}$。矩阵 $A$ 称为精度矩阵（precision matrix）。独立高斯观测的似然相乘时，精度相加，这往往比直接操作方差更方便。DDPM 的可计算后验正是一次这样的高斯乘法。

\begin{example}[一次线性高斯去噪]
\label{ex:diff-scalar-gaussian}
设 $X\sim\mathcal N(m,s^2)$，$Y=aX+b\varepsilon$，其中 $s,b>0$、$a\in\mathbb R$，噪声独立。则 $Y\sim\mathcal N(am,a^2s^2+b^2)$，且
\[
 \mathbb E[X\mid Y=y]=m+\frac{as^2}{a^2s^2+b^2}(y-am),\qquad
 \operatorname{Var}(X\mid Y=y)=\frac{s^2b^2}{a^2s^2+b^2}.
\]
例如 $m=0,s^2=4,a=1,b=1$ 时，观测 $y=3$ 对应条件均值 $2.4$、后验方差 $0.8$。估计比观测更靠近先验中心，但仍存在不确定性。把 $2.4$ 反复作为输出不能代表该条件高斯的抽样。
\end{example}

\subsection{梯度、雅可比与散度}
标量函数 $f:\mathbb R^d\to\mathbb R$ 的梯度 $\nabla f$ 是由各偏导数组成的列向量。向量场 $v:\mathbb R^d\to\mathbb R^d$ 的雅可比矩阵（Jacobian matrix）记作 $J_v$，元素为 $\partial v_i/\partial x_j$；散度（divergence）为 $\nabla\cdot v=\operatorname{tr}(J_v)$。梯度指向函数局部增长最快的方向，散度描述小体积中净流出的程度。二者作用不同：得分是对数密度的梯度，而流模型的密度变化由速度场的散度控制，不能把二者都含混称为“梯度”。

对正定高斯密度求对数并微分，得到 $\nabla_x\log\mathcal N(x;m,\Sigma)=-\Sigma^{-1}(x-m)$。这个向量在均值处为零，在远离均值时指向中心。若密度写成 $p(x)=\widetilde p(x)/Z$，归一化常数 $Z$ 不依赖 $x$，则 $\nabla_x\log p=\nabla_x\log\widetilde p$。这一消去常数的性质是得分建模的优势，但并不意味着知道梯度就能无需条件地恢复所有概率信息；不同连通分支间的相对权重、无穷远边界和归一化问题仍需考虑。

积分和微分交换需要条件。后文处理高斯卷积时，主要使用正噪声尺度下的平滑性，以及能把被积函数导数控制在可积函数之下的支配条件。分部积分则要求密度与相关向量场在无穷远足够快地衰减，或者明确边界条件。写出 $\nabla\int=\int\nabla$ 或丢弃边界项不是纯形式操作；对于尾部很重、支持有边界或噪声退化的问题，应重新检验成立范围。

\section{信息论、似然与变分推断}
\label{sec:diff-information}
\subsection{最大似然与相对熵}
给定独立训练样本，最大似然（maximum likelihood）寻找使 $\sum_i\log p_\theta(x_0^{(i)})$ 最大的参数。若总体密度存在且相关积分有限，最小化总体负对数似然等价于最小化真实分布相对模型分布的 Kullback--Leibler 散度，简称 KL 散度：
\begin{equation}
 D_{\mathrm{KL}}(q\|p)=\int q(x)\log\frac{q(x)}{p(x)}\,\mathrm dx.
 \label{eq:diff-kl}
\end{equation}
这里要求 $q$ 相对于 $p$ 绝对连续，否则散度可能为正无穷。KL 非负，当且仅当二者几乎处处相等时取零；它通常不对称，也不满足三角不等式，因此不是普通几何距离。

KL 的非负性可由 $\log u\leq u-1$ 证明。在分母非零并可积时，对 $u=p/q$ 应用不等式，再对 $q$ 求期望得到 $\mathbb E_q\log(p/q)\leq0$。从信息角度看，以错误分布描述样本会增加平均编码代价；从优化角度看，不同 KL 方向会对漏掉数据模式和在低数据区域分配质量产生不同惩罚。但“某个方向必然覆盖所有模式”不是无条件定理，因为模型族、约束、优化和有限样本都会限制最终结果。

连续变量的微分熵 $H(q)=-\int q\log q$ 可以为负，并随坐标缩放变化。负对数似然也可为负，不能仅凭其符号判定模型错误。若模型评价的是离散像素概率，每维比特数需要对离散观测概率取负对数，再除以 $d\log2$；若从连续密度换算，还要计入去量化和缩放项。扩散论文中的似然、变分上界、像素重建误差和感知指标衡量不同对象，不应互相替代。

\subsection{高斯 KL 的展开}
对于 $q=\mathcal N(m_q,S_q)$ 和 $p=\mathcal N(m_p,S_p)$，两协方差均正定，有
\begin{equation}
 D_{\mathrm{KL}}(q\|p)=\frac12\left[
 \operatorname{tr}(S_p^{-1}S_q)+(m_p-m_q)^\mathsf T S_p^{-1}(m_p-m_q)-d
 +\log\frac{\det S_p}{\det S_q}\right].
 \label{eq:diff-gaussian-kl}
\end{equation}
推导时展开两个高斯对数密度之差，再利用 $\mathbb E[(X-m)^\mathsf T A(X-m)]=\operatorname{tr}(A\operatorname{Cov}(X))$。若 $S_q=S_p=\sigma^2I_d$，矩阵迹、维数和行列式项抵消，只剩 $\|m_q-m_p\|^2/(2\sigma^2)$。因此，高斯条件分布之间的 KL 可以成为带权均方误差；它不是因为“平方误差比较常用”而被任意选择。

若模型方差也学习，则迹项与对数行列式项都依赖参数，不能只保留均值平方误差。增大预测方差会降低均值失配的惩罚，但同时提高对数行列式项，二者形成平衡。把方差固定是一项建模简化，通常方便优化和采样；把方差学习出来则需要保证正定性并正确处理相应损失。只训练噪声回归，并不能宣称已经最大化了一个任意可学习协方差模型的完整似然。

\subsection{潜变量、证据下界及其间隙}
若生成模型引入潜变量 $Z$，观测似然为 $p_\theta(x)=\int p_\theta(x,z)\,\mathrm dz$。积分往往难以直接计算。对一个满足适当支持条件的辅助分布 $r(z\mid x)$，定义证据下界（evidence lower bound, ELBO）
\begin{equation}
 \mathcal E(x)=\mathbb E_{r(z\mid x)}[\log p_\theta(x,Z)-\log r(Z\mid x)].
 \label{eq:diff-elbo-general}
\end{equation}
由贝叶斯公式拆分 $\log p_\theta(x,z)$，可精确得到
\[
 \log p_\theta(x)=\mathcal E(x)+D_{\mathrm{KL}}(r(z\mid x)\|p_\theta(z\mid x)).
\]
因此它是下界，且间隙就是辅助后验与真实模型后验的差异。也可用对数函数的凹性及 Jensen 不等式推导，但后验恒等式更清楚地解释了何时取等号。

在扩散中，潜变量是整条中间路径 $X_1,\ldots,X_T$，辅助分布是人为指定的前向扰动链。它通常固定且可抽样，不要求再训练一个复杂编码器。固定前向过程并不使 ELBO 自动等于真实对数似然；它只是让下界及其训练目标容易构造。与变分自编码器的联系在于潜变量和变分推断框架，与其区别在于潜变量组织成多层噪声路径、条件关系受到特定结构约束。高维中间状态也不等于每一层都压缩了向量维数。

\subsection{经验风险、随机梯度与总体结论的边界}
实际训练用样本平均近似总体期望，再用小批量近似样本平均。只要采样和权重正确、小批量估计具有所需可积性，其梯度可以是目标梯度的无偏估计；无偏不意味着方差小，也不保证非凸优化达到全局最优。对参数求导与取期望交换同样需要正则条件。常用可微网络与有限矩采样为这些操作提供合理基础，但极端权重和接近零的噪声尺度可能破坏数值稳定性。

损失值还受到归约方式影响：按全部维度求和、按每个像素平均、再按批量平均，得到的数值尺度不同。改变图像分辨率而继续使用总平方误差，会改变梯度量级；改为每维均值只改变一个常数，却可能要求不同学习率。比较实验必须说明损失定义和时间采样分布。数学上相差一个正常数的目标具有相同精确最优点，但有限步优化轨迹和停止准则未必相同。

\section{马尔可夫链与离散前向扩散}
\label{sec:diff-forward}
\subsection{局部随机转移如何规定整条路径}
随机序列具有马尔可夫性质（Markov property），是指给定当前状态后，下一状态的条件分布不再依赖更早历史。对前向扩散，这写为 $q(x_t\mid x_0,\ldots,x_{t-1})=q(x_t\mid x_{t-1})$，从而
\[
 q(x_{0:T})=p_{\mathrm{data}}(x_0)\prod_{t=1}^{T}q(x_t\mid x_{t-1}).
\]
若原始数据没有密度，上式的首项应理解为对 $P_{\mathrm{data}}$ 的积分，后续转移仍可具有密度。马尔可夫性质是条件独立，并非相邻状态无关；事实上噪声较小时，相邻图像高度相关。忽略这一区别会错误地把整条路径的分布写成各时刻边缘分布的乘积。

选择噪声日程（noise schedule）$0<\beta_t<1$，定义
\begin{equation}
 \alpha_t=1-\beta_t,\qquad \bar\alpha_t=\prod_{j=1}^{t}\alpha_j,\qquad \bar\alpha_0=1.
 \label{eq:diff-schedule}
\end{equation}
基本前向转移规定为
\begin{equation}
 X_t=\sqrt{\alpha_t}X_{t-1}+\sqrt{\beta_t}Z_t,
 \qquad Z_t\sim\mathcal N(0,I_d),
 \label{eq:diff-forward-step}
\end{equation}
其中各 $Z_t$ 相互独立并独立于 $X_0$。缩小信号并加入噪声的系数满足平方和为一，所以标准高斯是这个转移的不变分布。但任意输入的方差不必始终为一；若原始协方差为 $C_0$，第 $t$ 步协方差实际为 $\bar\alpha_t C_0+(1-\bar\alpha_t)I_d$。

\subsection{任意时间的闭式边缘}
\begin{proposition}[直接加噪公式]
\label{prop:diff-forward-marginal}
在式\eqref{eq:diff-forward-step}的独立性条件下，对任意 $1\leq t\leq T$，
\begin{equation}
 q(x_t\mid x_0)=\mathcal N(x_t;\sqrt{\bar\alpha_t}x_0,(1-\bar\alpha_t)I_d).
 \label{eq:diff-forward-marginal}
\end{equation}
因此可直接抽取独立的 $\varepsilon\sim\mathcal N(0,I_d)$，令 $X_t=\sqrt{\bar\alpha_t}X_0+\sqrt{1-\bar\alpha_t}\varepsilon$，获得该时刻的正确联合对 $(X_0,X_t)$。
\end{proposition}
\begin{proof}
逐步代入可知信号系数为 $\prod_{j=1}^{t}\sqrt{\alpha_j}=\sqrt{\bar\alpha_t}$，噪声是独立高斯的线性组合。记其每维方差为 $v_t$，则 $v_0=0$ 且 $v_t=\alpha_t v_{t-1}+\beta_t$。若 $v_{t-1}=1-\bar\alpha_{t-1}$，有 $v_t=\alpha_t(1-\bar\alpha_{t-1})+1-\alpha_t=1-\bar\alpha_t$，归纳完成。噪声与 $X_0$ 独立，所以这同时给出了条件均值和条件协方差。
\end{proof}

直接加噪使训练无需每次执行从零到 $t$ 的全部步骤。每个训练样本只需选择一个时间并构造该时刻观测，就能提供一个去噪任务。这里消除的是获取训练对的路径成本，生成时反向状态仍存在依赖，因此不能据此把所有反向步骤并行计算。若对多个时间使用同一个 $\varepsilon$，各时间的边缘仍正确，但它们之间的相关关系与原前向链通常不同；研究路径耦合、蒸馏或多时间损失时，必须明确到底使用哪种联合构造。

\subsection{信噪比与终点近似}
令 $a_t=\sqrt{\bar\alpha_t}$、$b_t=\sqrt{1-\bar\alpha_t}$。在信号每维方差标准化为一的参考尺度下，定义信噪比（signal-to-noise ratio, SNR）为 $\operatorname{SNR}_t=a_t^2/b_t^2$，对数信噪比为 $\lambda_t=\log(a_t^2/b_t^2)$。它衡量同一坐标尺度上信号与噪声的相对强弱，而非图像语义还能保留多少的精确测量。对于协方差差异显著的数据，不同方向的实际信噪比不同，单一标量只是描述扰动日程的便利记号。

若 $\bar\alpha_T$ 很小，终点信号被显著缩小，但 $q_T$ 并不在任意有限日程下精确等于标准高斯。可以给出一个清楚的充分估计：假设 $\mathbb E\|X_0\|^2<\infty$，把 $X_T=a_TX_0+b_T\varepsilon$ 与同一份 $\varepsilon$ 耦合，则
\[
 \mathbb E\|X_T-\varepsilon\|^2
 =a_T^2\mathbb E\|X_0\|^2+(b_T-1)^2d.
\]
它趋于零说明存在均方接近的耦合，也给出二阶 Wasserstein 距离的上界。这并不直接等同于 KL 或总变差距离的同样界，后两者需要另行分析。维数和数据尺度都出现在式中，因而“小残余系数”应结合数据范围判断。

这里的耦合（coupling）是把两个给定边缘分布放入同一个联合分布的方法。若 $P,Q$ 具有有限二阶矩，二阶 Wasserstein 距离定义为 $W_2(P,Q)=\inf_{(U,V)}(\mathbb E\|U-V\|^2)^{1/2}$，下确界遍及边缘分别为 $P,Q$ 的全部耦合。展示一个具体耦合只能给出上界，不保证它是最优耦合。总变差距离则比较所有事件概率差的上确界，关注的是另一种接近方式；不同距离的收敛不能在没有条件时互相替代。

日程可以通过指定每步方差，也可以通过指定累计信号 $\bar\alpha_t$ 来设计。后一种情况下由 $\beta_t=1-\bar\alpha_t/\bar\alpha_{t-1}$ 恢复转移，必须保证累计信号单调下降并处在有效范围内。不同日程把不同数量的训练时间分配给高噪声、低噪声和过渡区域。若时间均匀抽样，但对数信噪比变化极不均匀，实际学习任务并不是按噪声难度均匀分配；这为后续时间重加权提供了动机。

\section{反向条件分布及其精确可计算部分}
\label{sec:diff-posterior}
\subsection{已知干净样本时的后验}
训练阶段知道 $x_0$，因而能够计算 $q(x_{t-1}\mid x_t,x_0)$。对 $t\geq2$，贝叶斯公式及马尔可夫性质给出该分布正比于 $q(x_t\mid x_{t-1})q(x_{t-1}\mid x_0)$。令待求变量为 $u=x_{t-1}$，把两个高斯指数相加，其中与 $u$ 有关的部分为
\[
 -\frac{\|x_t-\sqrt{\alpha_t}u\|^2}{2\beta_t}
 -\frac{\|u-\sqrt{\bar\alpha_{t-1}}x_0\|^2}{2(1-\bar\alpha_{t-1})}.
\]
二次项系数为 $A_t=\alpha_t/\beta_t+1/(1-\bar\alpha_{t-1})$，一次项向量为 $h_t=\sqrt{\alpha_t}x_t/\beta_t+\sqrt{\bar\alpha_{t-1}}x_0/(1-\bar\alpha_{t-1})$。利用高斯配方，方差是 $A_t^{-1}$，均值是 $A_t^{-1}h_t$。

化简时使用 $\bar\alpha_t=\alpha_t\bar\alpha_{t-1}$，得到
\begin{align}
 \widetilde\beta_t&=\frac{1-\bar\alpha_{t-1}}{1-\bar\alpha_t}\beta_t,\label{eq:diff-posterior-var}\\
 \widetilde\mu_t(x_t,x_0)&=
 \frac{\sqrt{\bar\alpha_{t-1}}\beta_t}{1-\bar\alpha_t}x_0+
 \frac{\sqrt{\alpha_t}(1-\bar\alpha_{t-1})}{1-\bar\alpha_t}x_t,
 \label{eq:diff-posterior-mean}\\
 q(x_{t-1}\mid x_t,x_0)&=\mathcal N(x_{t-1};\widetilde\mu_t,\widetilde\beta_t I_d).
 \label{eq:diff-posterior}
\end{align}
由于 $t\geq2$ 时方差为正，这是真正的高斯密度。形式上令 $t=1$，有 $\widetilde\beta_1=0$，因为给定 $x_0$ 后 $X_0$ 没有任何随机性；此时条件分布是点质量，不能把零方差直接代入普通高斯 KL。第一步的观测模型必须在损失分解中单独处理。

\begin{example}[两步扩散的数值核对]
取 $\beta_1=0.1,\beta_2=0.2$，则 $\alpha_1=0.9,\alpha_2=0.8,\bar\alpha_2=0.72$。在 $t=2$，后验方差为 $0.2\times0.1/0.28=1/14$，后验均值为 $0.67763x_0+0.31944x_2$，系数保留五位小数。它不是两个观测的普通凸组合，因为前向状态经过了尺度收缩。若 $x_0=1,x_2=0.5$，均值约为 $0.83735$，同时仍需保留方差所描述的不确定性。
\end{example}

\subsection{为什么真正的反向转移通常不是高斯}
生成阶段只知道 $x_t$，不知道 $x_0$。正确的反向分布应为
\begin{equation}
 q(x_{t-1}\mid x_t)=\int q(x_{t-1}\mid x_t,x_0)q(x_0\mid x_t)\,\mathrm dx_0.
 \label{eq:diff-reverse-mixture}
\end{equation}
若数据含多个模式，右侧是许多不同均值高斯的混合，通常并不是单一高斯。所谓“前向每步是高斯，所以反向每步也是高斯”遗漏了对未知数据的边缘化。有限步 DDPM 使用条件高斯来近似反向转移是一项模型选择；当时间步很小时，在适当平滑条件下可由连续扩散理论支持局部近似，但不能推成任意大跨度反向跳跃都精确高斯。

把式\eqref{eq:diff-posterior-mean}中 $x_0$ 的系数记为 $c_t$，$x_t$ 的系数记为 $d_t$。全期望与全协方差公式给出
\begin{align}
 \mathbb E[X_{t-1}\mid X_t]&=c_t\mathbb E[X_0\mid X_t]+d_tX_t,\label{eq:diff-reverse-mean}\\
 \operatorname{Cov}(X_{t-1}\mid X_t)&=\widetilde\beta_t I_d+c_t^2\operatorname{Cov}(X_0\mid X_t).
 \label{eq:diff-reverse-var}
\end{align}
第二项来自无法确定原始样本的额外不确定性。因此把反向方差固定为 $\widetilde\beta_t$，一般不能恢复完整真实反向协方差。它在建模中仍可能有用，但必须区分方便的方差设置与精确条件分布。这个区别也解释了为什么学习方差、缩短单步跨度或改用其他采样动力学可能产生不同结果。

\subsection{从可计算教师信号到可部署模型}
通常设 $p_\theta(x_{t-1}\mid x_t)=\mathcal N(x_{t-1};\mu_\theta(x_t,t),\gamma_t^2I_d)$，其中 $\gamma_t>0$ 为选定的反向标准差。符号 $\gamma_t$ 与前向噪声 $b_t$ 分开，避免把两种不同方差混淆。网络以当前状态和时间为输入，拟合均值或与均值等价的量。训练时可利用 $x_0$ 构造监督，并不代表网络在测试时也能访问 $x_0$；监督变量多于输入变量，是一般回归问题的常见结构。

若直接对 $\widetilde\mu_t(X_t,X_0)$ 做平方回归，命题\ref{prop:diff-regression}说明最优网络输出为 $\mathbb E[\widetilde\mu_t(X_t,X_0)\mid X_t]$，即真实反向条件均值。后验均值标签本身依赖每个训练样本，但通过许多样本平均可以学习一个只依赖可观测输入的函数。噪声预测只是对这个回归问题的一种重新参数化，它并没有绕过条件期望，也不要求从单个受扰观测唯一恢复隐藏的全部信息。

如果原始分布恰为标准高斯，前向过程所有边缘都是标准高斯，真正的反向转移也可精确写为 $\mathcal N(\sqrt{\alpha_t}x_t,\beta_t I_d)$。此时真实反向方差是 $\beta_t$，而不是通常更小的 $\widetilde\beta_t$。这一完全可解的特例非常适合检查系数、方差和时间方向：它说明即便网络达到理想预测，也要为“精确采样”配置正确转移，而不能只检查均值是否合理。

\section{DDPM 的变分目标与噪声回归}
\label{sec:diff-ddpm-loss}
\subsection{反向联合模型和路径下界}
选择终点先验 $p_T=\mathcal N(0,I_d)$，并以反向转移规定
\begin{equation}
 p_\theta(x_{0:T})=p_T(x_T)\prod_{t=1}^{T}p_\theta(x_{t-1}\mid x_t).
 \label{eq:diff-generative-joint}
\end{equation}
观测分布是对全部中间状态积分后的 $p_\theta(x_0)$。把固定的 $q(x_{1:T}\mid x_0)$ 作为辅助后验，式\eqref{eq:diff-elbo-general}直接给出负对数似然上界
\[
 -\log p_\theta(x_0)\leq
 \mathbb E_q\left[\log\frac{q(X_{1:T}\mid x_0)}{p_\theta(x_0,X_{1:T})}\right].
\]
右侧的期望只需要从已知前向过程抽样。它不是把生成链在训练时完整运行一遍，再与原图比较，而是利用一个便于计算的辅助分布来组织各局部学习问题。

对给定 $x_0$ 的前向路径使用反向链式分解，有
\[
 q(x_{1:T}\mid x_0)=q(x_T\mid x_0)\prod_{t=2}^{T}q(x_{t-1}\mid x_t,x_0).
\]
该等式依赖前向链的马尔可夫结构，可由连续应用贝叶斯公式验证；分子与分母中的中间边缘项会逐项抵消。代入上界并按对应变量取期望，得到
\begin{align}
 \mathcal L_{\mathrm{VLB}}(x_0)
 &=D_{\mathrm{KL}}(q(X_T\mid x_0)\|p_T)\nonumber\\
 &\quad+\sum_{t=2}^{T}\mathbb E_{q(X_t\mid x_0)}
 D_{\mathrm{KL}}(q(X_{t-1}\mid X_t,x_0)\|p_\theta(X_{t-1}\mid X_t))\nonumber\\
 &\quad-\mathbb E_{q(X_1\mid x_0)}\log p_\theta(x_0\mid X_1).
 \label{eq:diff-vlb}
\end{align}
三部分依次约束终点与先验的匹配、中间反向转移的匹配以及最后一步的观测恢复。若前向日程与先验固定，第一项对网络参数是常数；这只意味着它不提供训练梯度，不意味着其数值为零或可以在报告似然界时删除。

\subsection{从均值误差到噪声误差}
当反向方差 $\gamma_t^2I_d$ 固定时，高斯 KL 中关于参数的部分是 $\|\widetilde\mu_t-\mu_\theta\|^2/(2\gamma_t^2)$。利用 $x_t=a_tx_0+b_t\varepsilon$ 消去 $x_0$，式\eqref{eq:diff-posterior-mean}可改写为
\[
 \widetilde\mu_t=\frac1{\sqrt{\alpha_t}}\left(x_t-\frac{\beta_t}{b_t}\varepsilon\right).
\]
于是令网络预测噪声，并定义
\begin{equation}
 \mu_\theta(x_t,t)=\frac1{\sqrt{\alpha_t}}\left(x_t-\frac{\beta_t}{b_t}\varepsilon_\theta(x_t,t)\right),
 \label{eq:diff-eps-mean}
\end{equation}
中间 KL 的可训练部分变成
\begin{equation}
 w_t\mathbb E\|\varepsilon-\varepsilon_\theta(a_tX_0+b_t\varepsilon,t)\|^2,
 \qquad w_t=\frac{\beta_t^2}{2\gamma_t^2\alpha_t(1-\bar\alpha_t)}.
 \label{eq:diff-vlb-weight}
\end{equation}
系数中的每一项都有来源：$\beta_t/b_t$ 来自均值参数化，$\alpha_t$ 来自收缩因子，$\gamma_t^2$ 来自高斯精度。漏掉这些系数并不一定使训练不可用，却会改变目标与原变分界的关系。

常见的简化目标取时间均匀分布并删除显式的 $w_t$，写为
\begin{equation}
 \mathcal L_{\mathrm{simple}}=
 \mathbb E_{t\sim\operatorname{Unif}\{1,\ldots,T\},X_0,\varepsilon}
 \|\varepsilon-\varepsilon_\theta(a_tX_0+b_t\varepsilon,t)\|^2.
 \label{eq:diff-simple}
\end{equation}
这是一种重新加权的去噪回归目标，不等于逐项不变地优化完整变分下界。原始 DDPM 工作把此类参数化与训练选择联系起来\cite{ho2020}；本节给出的等式可以独立检验。对无限容量函数，正时间权重不改变每个时间的条件均值解；对共享网络和有限训练预算，改变权重会改变实际优化重点。

\subsection{时间抽样、第一步和方差学习}
若目标是 $\sum_{t=1}^{T}A_t$，却按概率 $\pi_t>0$ 抽时间，则 $A_t/\pi_t$ 的单次抽样期望才等于原和；若目标是平均值，再除以 $T$。不加补偿地改变时间分布，相当于改变目标。重要性抽样（importance sampling）可以降低估计方差，但极小 $\pi_t$ 会产生大权重，反而造成不稳定。选择抽样分布时应平衡各时间损失或梯度的波动，而不是只优先选择当前损失最大的时间。

最后的 $p_\theta(x_0\mid x_1)$ 需要匹配数据表示。连续数据可以采用具有正方差的观测密度；离散像素可以对相应区间积分得到概率质量。实践中最后一步有时直接输出均值，这是一种输出或采样约定，不应与训练时定义的完整观测概率混为一谈。若把 $t=1$ 简单纳入噪声回归，这同样是一种训练设计；要把所得数值报告为似然界，仍须恢复正确的重建项和常数。

学习反向方差时，可以用受约束的参数化保证其为正，再结合完整 KL 或混合目标训练。Improved DDPM 讨论了学习方差和其他训练改进\cite{nichol2021}。理解这类方法时，应检查均值与方差分别接受哪些梯度、目标是否包含停止梯度的设计，以及采样器是否使用学到的方差。较低的噪声均方误差不自动表示方差校准更好；若关注概率似然，两部分都属于模型。

\section{预测参数化、信噪比权重与训练含义}
\label{sec:diff-parameterization}
\subsection{预测噪声与预测干净数据}
只要 $a_t,b_t>0$，噪声预测与干净样本预测之间具有逐点代数关系
\begin{equation}
 \widehat x_{0,\theta}=\frac{x_t-b_t\varepsilon_\theta}{a_t},\qquad
 \widehat\varepsilon_\theta=\frac{x_t-a_tD_\theta}{b_t},
 \label{eq:diff-x-eps-convert}
\end{equation}
其中 $D_\theta$ 是直接预测干净数据的网络输出。通过一致变换表示同一个函数时，两种参数化不改变可表示的理想估计；但把同一结构的网络分别训练成二者，并使用相同的无权均方误差，通常会得到不同优化问题。因为
\[
 \|\varepsilon-\widehat\varepsilon_\theta\|^2
 =\frac{a_t^2}{b_t^2}\|X_0-D_\theta\|^2
 =\operatorname{SNR}_t\|X_0-D_\theta\|^2.
\]
无权噪声预测对应按信噪比加权的干净数据误差；不能同时声称两者损失完全相同而又忽略这个时间权重。

在高噪声端 $a_t$ 很小，从噪声输出反求 $x_0$ 需要除以小数，噪声预测的微小误差可能被放大。在低噪声端 $b_t$ 很小，从干净数据输出反求噪声也有类似问题。这是参数化的数值条件问题，不表示对应概率对象不存在。端点处若某个系数为零，代数换算失效，应使用端点边界条件或直接输出定义；推导时把内部区间有效公式不加说明地延伸到端点，容易制造人为奇异性。

\subsection{速度参数化及其正交结构}
在 $a_t^2+b_t^2=1$ 的方差保持设置中，可定义速度目标（velocity parameterization）
\begin{equation}
 V_t=a_t\varepsilon-b_tX_0.
 \label{eq:diff-v-target}
\end{equation}
它与 $X_t$ 共同组成一个正交变换：
\[
 \begin{pmatrix}X_t\\V_t\end{pmatrix}
 =\begin{pmatrix}a_t&b_t\\-b_t&a_t\end{pmatrix}
 \begin{pmatrix}X_0\\\varepsilon\end{pmatrix},\qquad
 X_0=a_tX_t-b_tV_t,\quad\varepsilon=b_tX_t+a_tV_t.
\]
因此从网络 $v_\theta$ 恢复 $x_0$ 或噪声不需要分别除以 $a_t$ 或 $b_t$。若写 $a=\cos\varphi,b=\sin\varphi$，则同一对 $(X_0,\varepsilon)$ 形成的插值对角度的导数恰为 $V$，这说明“速度”一词的来源。它并不是任意时间坐标下无须缩放的实际轨迹速度，更不是流匹配所有文献中速度符号的统一定义。

若干净数据中心化且每维方差为一，独立性使 $V_t$ 的协方差也为单位矩阵，提供了目标尺度相对稳定的动机。若真实数据方差不同，这一结论随之改变。对于由同一个 $v_\theta$ 转换出的预测，有 $\|X_0-\widehat x_0\|^2=b_t^2\|V_t-v_\theta\|^2$，且 $\|\varepsilon-\widehat\varepsilon\|^2=a_t^2\|V_t-v_\theta\|^2$。这进一步表明参数化和损失权重是两个需要联合考虑的选择。

\subsection{去噪器的预条件与尺度设计}
神经网络对输入输出尺度敏感。可把去噪函数表示为
\[
 D_\theta(y,\sigma)=c_{\mathrm{skip}}(\sigma)y+
 c_{\mathrm{out}}(\sigma)F_\theta(c_{\mathrm{in}}(\sigma)y,c_{\mathrm{noise}}(\sigma)),
\]
其中 $\sigma$ 是加性噪声标准差，四个 $c$ 是事先设计的标量函数，$F_\theta$ 为主要网络。跳跃项承担容易预测的部分，输入缩放控制激活幅度，输出缩放控制残差量级，噪声嵌入则告诉网络当前难度。这种预条件（preconditioning）改变了网络需要拟合的数值形状，不应与对数据概率分布实施额外物理过程混淆。EDM 对这类设计及采样选择作了系统讨论\cite{karras2022}。

以加性高斯观测 $Y=X+\sigma\varepsilon$ 为例，若先验近似各向同性高斯、每维方差为 $\sigma_{\mathrm{data}}^2$，线性最优收缩系数为 $\sigma_{\mathrm{data}}^2/(\sigma_{\mathrm{data}}^2+\sigma^2)$。把类似系数放入跳跃支路，可让网络专注于偏离简单高斯假设的非线性结构。这里的高斯先验用于解释预条件的动机，并不要求真实图像服从单峰高斯，也不保证这些系数在所有模态上最优。

\subsection{自条件、预测误差与生成误差}
自条件（self-conditioning）把先前得到的干净样本估计作为额外输入，帮助网络迭代修正表示。它不等于把真实 $x_0$ 提供给生成阶段，也不是对未知答案的合法性豁免；训练时提供的自条件信息必须来自允许的估计路径，并与推理过程相容。若训练总使用近乎完美的辅助估计，推理却使用有偏估计，会形成新的分布偏移。理解这类扩展应先确定额外输入的来源、时间和梯度路径，再讨论收益。

训练回归误差是在真实扰动分布上平均的，而采样过程中网络接收到的状态来自近似生成轨迹。两者不完全一致，小误差可能经多步动力学累积，也可能被稳定结构部分抑制。误差放大与向量场的局部 Lipschitz 常数、步长、引导尺度和端点行为有关。因此更低的平均训练损失不保证每一种采样器下都获得更好的样本。判断学习是否有效，应同时查看不同噪声区间的预测、生成分布、条件满足情况以及计算成本。

\section{得分函数、去噪恒等式与得分匹配}
\label{sec:diff-score}
\subsection{得分描述什么，不能单独描述什么}
对正且可微的密度 $p_t(x)$，关于数据变量的得分函数（score function）定义为 $s_t(x)=\nabla_x\log p_t(x)$。统计学也把参数导数 $\nabla_\theta\log p_\theta(x)$ 称为得分，本文使用的是前者，必须注意求导变量。$s_t(x)$ 给出局部增加对数密度的方向和速率，但不直接给出该位置的概率质量。一个高维分布可以在中心具有最大密度，同时绝大部分质量集中在离中心较远的壳层上，因此总向密度最大点移动并不等于正确采样。

标准高斯是最简单的例子：$s(x)=-x$。若只沿得分做梯度上升，所有点都会趋向零，最终分布成为点质量；若希望维持标准高斯，需要合适的随机扰动或不同的输运动力学。对于多峰分布，得分在每个局部峰附近指向该峰，低密度区域的估计又较难，这使单一尺度的采样容易困在一个区域。多尺度加噪能让分离的模式在高噪声下逐渐连通，为跨模式移动提供较平滑的中间问题，但具体混合速度仍依赖几何和算法。

\subsection{条件得分的平均就是边缘得分}
令 $Y=aX+b\varepsilon$，其中 $a>0,b>0$，$\varepsilon\sim\mathcal N(0,I_d)$ 独立于 $X$，并假定二阶矩有限。记平滑边缘密度为 $p_Y$。在允许把微分移入积分的条件下，
\begin{align}
 \nabla_y p_Y(y)
 &=\int q(y\mid x)\nabla_y\log q(y\mid x)\,\mathrm dP_X(x),\nonumber\\
 \nabla_y\log p_Y(y)
 &=\mathbb E\!\left[-\frac{Y-aX}{b^2}\,\middle|\,Y=y\right]
 =-\frac1b\mathbb E[\varepsilon\mid Y=y].
 \label{eq:diff-score-identity}
\end{align}
第一行只是在高斯核上求导，第二行则除以正密度并识别出后验权重。原始分布可含点质量，积分使用概率测度即可。这是边缘化后求得分的准确方法，不能直接把某个训练样本的条件得分当成全部数据的得分。

整理式\eqref{eq:diff-score-identity}得到去噪恒等式，也常称为高斯噪声下的 Tweedie 公式：
\begin{equation}
 \mathbb E[X\mid Y=y]=\frac{y+b^2\nabla_y\log p_Y(y)}{a}.
 \label{eq:diff-tweedie}
\end{equation}
因此准确去噪器、准确条件噪声预测和准确平滑密度得分包含等价信息。以 $a=1$ 为例，$b^2s(y)$ 是从观测走向最优均方去噪结果的修正量。公式并没有说沿这个方向移动一次就得到后验样本，也没有说任意神经网络预测都必然对应某个全局归一化密度的精确梯度。

\subsection{从不可观测得分到可计算监督}
想直接最小化 $\mathbb E\|s_\theta(Y)-\nabla\log p_Y(Y)\|^2$，困难在于真实边缘密度未知。可改为去噪得分匹配（denoising score matching）目标
\begin{equation}
 \mathcal J_b(\theta)=\mathbb E\left\|s_\theta(aX+b\varepsilon)+\frac{\varepsilon}{b}\right\|^2.
 \label{eq:diff-dsm}
\end{equation}
由命题\ref{prop:diff-regression}与式\eqref{eq:diff-score-identity}，它等于真实得分误差加上一个与网络无关的条件方差项。这个证明只需要平方可积性和条件期望，不需要知道边缘密度的归一化常数。训练标签 $-\varepsilon/b$ 在每次实验中不同，但其给定观测后的平均恰好是需要学习的函数。

代入 $s_\theta(x_t,t)=-\varepsilon_\theta(x_t,t)/b_t$，有 $\|s_\theta+\varepsilon/b_t\|^2=\|\varepsilon_\theta-\varepsilon\|^2/b_t^2$。所以无权噪声预测对应以 $b_t^2$ 加权的去噪得分误差。这条关系把 DDPM 的回归目标与多尺度得分学习联系起来，也说明比较不同目标时必须携带噪声尺度。低噪声时条件得分标签的幅度可能很大，权重和预条件对梯度稳定性尤其重要。

经典得分匹配还有另一种形式。假设 $p$ 光滑，相关二阶量可积，边界项消失，则分部积分给出
\[
 \frac12\mathbb E_p\|s_\theta-\nabla\log p\|^2
 =\mathbb E_p\left[\frac12\|s_\theta\|^2+\nabla\cdot s_\theta\right]+C,
\]
其中 $C$ 与参数无关。其依据是 $\int s_\theta\cdot\nabla p=-\int p\nabla\cdot s_\theta$。这一框架的基础工作见\cite{hyvarinen2005}。高维神经网络中散度计算可能昂贵；去噪形式利用已知扰动核，直接提供回归标签，因而在扩散学习中非常方便。

\subsection{朗之万动力学与噪声退火}
对于固定目标密度 $p$，在适当边界和遍历条件下，朗之万动力学（Langevin dynamics）$\mathrm dY_u=\nabla\log p(Y_u)\,\mathrm du+\sqrt2\,\mathrm dW_u$ 以 $p$ 为不变密度。对应的有限步近似为 $Y_{k+1}=Y_k+h\nabla\log p(Y_k)+\sqrt{2h}\,Z_k$。梯度项把状态拉向较高密度，随机项补充扩散；二者系数的配合很重要。去掉随机项会变成寻找模式，任意改变噪声大小则一般改变平稳分布。

有限步长近似通常有离散化偏差，并非迭代一次就得到目标分布；多峰情况下混合可能很慢。退火朗之万方法先在高噪声平滑分布上移动，再逐步降低噪声，利用前一尺度结果初始化下一尺度。它与沿时间反转一条指定扩散路径有关联，但不应把每个尺度上的平衡采样和一次反向扩散步完全等同。两者使用的动力学、时间参数和误差来源不同，是否需要在每个尺度充分混合也应按具体方法判断。

\section{连续时间前置知识：布朗运动与随机微分方程}
\label{sec:diff-sde}
\subsection{从小步加噪到连续模型}
令离散步长为 $\Delta\tau$，取 $\beta_t\approx\beta(\tau)\Delta\tau$。利用 $\sqrt{1-u}=1-u/2+O(u^2)$，前向增量近似为 $-\tfrac12\beta(\tau)X\Delta\tau+\sqrt{\beta(\tau)\Delta\tau}\,Z$。确定部分按时间步长缩放，随机部分按其平方根缩放，因此随机波动不能被当作普通可微函数的速度。连续极限自然引出随机微分方程（stochastic differential equation, SDE），其中随机项由布朗运动驱动。

$d$ 维标准布朗运动（Brownian motion）$W_\tau$ 从零出发，路径连续，互不相交时间区间的增量独立，且 $W_{\tau+h}-W_\tau\sim\mathcal N(0,hI_d)$。典型路径不可按普通微积分求导，所谓白噪声只是形式上的导数记号。理解 $\mathrm dW$ 时，应记住其短时间方差是 $\mathrm d\tau$，而不是把它看成一个可随意移项约分的普通小量。随机积分的定义决定后续微分法则，本文采用伊藤（Itô）约定。

考虑空间无关的各向同性扩散系数 $g(\tau)$，方程为
\begin{equation}
 \mathrm dX_\tau=f(X_\tau,\tau)\,\mathrm d\tau+g(\tau)\,\mathrm dW_\tau.
 \label{eq:diff-forward-sde}
\end{equation}
其严格含义是积分关系 $X_\tau=X_0+\int_0^\tau f(X_r,r)\,\mathrm dr+\int_0^\tau g(r)\,\mathrm dW_r$。对漂移在空间中全局 Lipschitz、至多线性增长，且扩散系数满足相应连续性和可积性的一类情况，可引用存在唯一强解的标准定理。实际神经场或奇异端点可能不满足这组简便的充分条件，应在有限时间区间内另行控制，而不能无条件宣称解必然良好。

全局 Lipschitz 条件可写为 $\|f(x,\tau)-f(y,\tau)\|\leq L\|x-y\|$，其中有限常数 $L$ 对所考虑时间区间统一有效；线性增长可写为 $\|f(x,\tau)\|\leq C(1+\|x\|)$。它们分别限制对状态扰动的敏感程度和无穷远增长。强解表示在已给定的概率空间、初始状态和布朗运动上构造适应的解过程；弱解允许连同概率空间和驱动噪声一起构造。这里的“强、弱”是随机方程术语，与网络能力高低无关，也不等于数值算法误差的大小。

\subsection{伊藤修正与生成算子}
若测试函数 $\phi(x,\tau)$ 对空间二阶连续可微、对时间一阶连续可微，并满足使积分有意义的增长条件，伊藤公式给出
\begin{equation}
 \mathrm d\phi(X_\tau,\tau)=
 \left(\partial_\tau\phi+f\cdot\nabla\phi+\tfrac12g^2\Delta\phi\right)\mathrm d\tau
 +g\nabla\phi^\mathsf T\mathrm dW_\tau.
 \label{eq:diff-ito}
\end{equation}
$\Delta\phi=\sum_j\partial_{x_j}^2\phi$ 是拉普拉斯算子。多出的二阶项源自布朗增量的二次变差，不能用普通链式法则删掉。例如一维 $\phi(x)=x^2$、$X=W$ 时，$\mathrm d(W^2)=2W\,\mathrm dW+\mathrm d\tau$，取期望得到 $\mathbb E[W_\tau^2]=\tau$，与增量定义一致。

算子 $\mathcal A_\tau\phi=f\cdot\nabla\phi+\tfrac12g^2\Delta\phi$ 称为无显式时间导数部分的生成算子（generator）。在随机积分均值为零的条件下，$\mathrm d\mathbb E[\phi(X_\tau)]/\mathrm d\tau=\mathbb E[\mathcal A_\tau\phi(X_\tau)]$。它把随机路径问题转化为任意测试函数期望的演化，再通过分部积分得到密度方程。后文的反向与概率流关系都依赖这一步，因此仅把 SDE 看作“普通微分方程加一点噪声”是不够的。

\subsection{福克尔--普朗克方程}
若 $X_\tau$ 存在足够光滑的密度 $p_\tau$，并允许对紧支撑光滑测试函数使用上述恒等式，则密度满足福克尔--普朗克方程（Fokker--Planck equation）
\begin{equation}
 \partial_\tau p_\tau=-\nabla\cdot(fp_\tau)+\frac12g(\tau)^2\Delta p_\tau.
 \label{eq:diff-fp}
\end{equation}
这里空间无关的 $g$ 才能直接移出空间导数。漂移项搬运概率质量，扩散项让局部分布平滑。积分后总概率保持为一，需要无穷远通量消失或适当的边界条件。方程可先在弱意义下理解；若密度足够正则，再转为逐点成立的偏微分方程。

对状态依赖的矩阵扩散系数，扩散项应写成协方差矩阵与密度乘积的二阶偏导和，反向漂移也会出现相应空间导数项。本文主线只处理 $g(\tau)I_d$，是为了让关键关系清楚可核查；不能把这里的反向公式未经修改用于任意乘性噪声或受约束流形。数据位于球面、旋转群或其他空间时，体积测度、梯度和扩散算子都需要与几何相容。

\subsection{方差保持与方差增长两类过程}
取 $f(x,\tau)=-\tfrac12\beta(\tau)x$、$g(\tau)=\sqrt{\beta(\tau)}$，得到方差保持（variance preserving, VP）SDE。设 $B(\tau)=\int_0^\tau\beta(r)\,\mathrm dr$，则条件分布为
\[
 X_\tau\mid X_0=x_0\sim\mathcal N(e^{-B(\tau)/2}x_0,[1-e^{-B(\tau)}]I_d).
\]
其推导使用线性方程的积分因子以及伊藤积分的等距性质。名称中的方差保持，准确地说是标准高斯输入在该过程下保持不变，而不是任意数据的实际方差恒定不变。它与离散 DDPM 的累计系数具有直接对应。

若取零漂移，令非减可微函数 $v(\tau)$ 满足 $v(0)=0$，并设 $g(\tau)=\sqrt{v'(\tau)}$，则 $X_\tau=X_0+\sqrt{v(\tau)}\varepsilon$ 在边缘意义上成立。数据不收缩，噪声方差不断增加，这提供方差增长（variance exploding, VE）过程的一种基本表述。若实际采用正的最小噪声尺度，则起点已经是平滑数据分布，应明确其与真实数据之间的差异。VP 与 VE 可以通过坐标缩放和时间重参数化建立联系，但网络尺度、训练权重和采样步长也必须一同换算。

\section{时间反演、概率流与连续密度计算}
\label{sec:diff-reversal}
\subsection{反向 SDE 的时间符号}
设式\eqref{eq:diff-forward-sde}在所研究时间区间内具有正且足够光滑的边缘密度 $p_\tau$，系数与密度满足反向扩散定理所需的正则性、非爆炸性及积分条件。对于空间无关的 $g$，反向时间表示为
\begin{equation}
 \mathrm dX_\tau=[f(X_\tau,\tau)-g(\tau)^2\nabla\log p_\tau(X_\tau)]\,\mathrm d\tau
 +g(\tau)\,\mathrm d\overline W_\tau,
 \label{eq:diff-reverse-sde}
\end{equation}
其中时间从 $\mathcal T$ 向零递减，反向布朗增量的协方差按 $|\mathrm d\tau|$ 计算。反向过程定理是随机分析结论，本文引用其在得分扩散框架中的表述\cite{song2021sde}；下文通过密度关系和可解例子解释符号，但不把形式计算当成完整的路径测度证明。

为避免负时间增量造成混淆，设递增反向时钟 $u=\mathcal T-\tau$，$Y_u=X_{\mathcal T-u}$。则等价写法为
\[
 \mathrm dY_u=[-f(Y_u,\mathcal T-u)+g(\mathcal T-u)^2s_{\mathcal T-u}(Y_u)]\,\mathrm du
 +g(\mathcal T-u)\,\mathrm dB_u,
\]
其中 $B_u$ 是相应反向滤过下的标准布朗运动。这样可直接看到得分项沿增加密度的方向起作用。把反向公式中的负号与正步长混用，会把去噪变成进一步扰乱。核查公式时应先写清时钟递增还是递减，再判断漂移的方向。

反向过程仍包含随机噪声，因为它在每一步抽取一个可能的前驱状态，而不是恢复唯一历史。当前观测通常对应多种过去，正确反演必须分配这些可能性。如果初始使用的噪声先验与真正的 $p_{\mathcal T}$ 不同，即使得分完全准确，也未必精确回到原数据分布；这属于终点失配。若得分近似、积分步长有限，还分别增加统计误差和离散化误差。三种误差应分开讨论，不能统一归咎于“去噪不够强”。

\subsection{从密度守恒推导概率流 ODE}
确定性常微分方程（ordinary differential equation, ODE）$\mathrm dX_\tau/\mathrm d\tau=b(X_\tau,\tau)$ 诱导的密度满足连续性方程（continuity equation）$\partial_\tau p=-\nabla\cdot(bp)$。由于 $p_\tau\nabla\log p_\tau=\nabla p_\tau$，选择
\begin{equation}
 b(x,\tau)=f(x,\tau)-\frac12g(\tau)^2\nabla\log p_\tau(x)
 \label{eq:diff-pf-field}
\end{equation}
时，其连续性方程右侧恰为式\eqref{eq:diff-fp}。在密度方程和流的解具有所需存在唯一性、且初始边缘一致的条件下，这个概率流 ODE（probability flow ODE）与前向 SDE 具有相同的每时刻边缘密度。反向积分它即可由终点分布返回数据分布。

共享边缘分布并不等于共享路径分布。ODE 在给定起点后确定地沿一条轨迹运动，SDE 即便起点相同也存在随机路径；二者在每个单独时刻的样本群体可以服从相同分布。ODE 的生成随机性来自随机初始状态，因此确定性采样器并不意味着所有生成样本相同。应把“固定种子后的可重复性”与“总体分布是否随机”区分开，这一点也适用于普通正规化流。

将真实得分替换为网络 $s_\theta$ 后，两种动力学一般不再自动具有相同边缘。因为概率流推导使用的是正在演化的真实密度的得分，而近似网络未必等于任一近似过程当前密度的得分。只有在估计准确、初始分布匹配、积分精确等理想条件下，前述等价结论才能完整转移。现实中可以把 SDE 与 ODE 看成利用同一学习信息的不同采样方式，其表现差异需要实际评价。

\begin{example}[标准高斯的时间方向检查]
令 VP 过程从标准高斯出发，则 $p_\tau=\mathcal N(0,I_d)$、$s_\tau(x)=-x$。概率流速度为 $-\tfrac12\beta x-\tfrac12\beta(-x)=0$，所以每个 ODE 状态保持不变，群体分布当然也不变。反向递增时钟下的 SDE 漂移为 $\tfrac12\beta x+\beta(-x)=-\tfrac12\beta x$，仍是以标准高斯为不变分布的收缩扩散。这同时检验了得分符号和概率流中的二分之一系数。
\end{example}

\subsection{沿流的对数密度与似然}
若 $b$ 对空间可微并生成可逆流，密度为正，则沿轨迹有
\begin{equation}
 \frac{\mathrm d}{\mathrm d\tau}\log p_\tau(X_\tau)=-\nabla\cdot b(X_\tau,\tau).
 \label{eq:diff-density-ode}
\end{equation}
证明只需把连续性方程除以 $p$，再加上链式法则中的 $b\cdot\nabla\log p$，两个输运项抵消。因此正向把数据点送到终点时，$\log p_0(x_0)=\log p_{\mathcal T}(x_{\mathcal T})+\int_0^{\mathcal T}\nabla\cdot b(X_\tau,\tau)\,\mathrm d\tau$。这是连续正规化流的密度换元关系，并不是从回归损失数值直接读出似然。

高维散度可用随机迹估计减少计算。对独立辅助向量 $R$，若 $\mathbb E[RR^\mathsf T]=I_d$，则给定雅可比 $J_b$ 有 $\mathbb E_R[R^\mathsf T J_bR]=\operatorname{tr}(J_b)$。这只是迹估计的无偏性；有限探针引入随机误差，ODE 求解引入数值误差，模型本身还可能不准确。所以“允许似然计算”与“无需误差地获得真实数据似然”含义不同。若原始数据分布奇异，通常在正的最小噪声或明确去量化模型下讨论密度。

\section{采样原理：祖先抽样、DDIM 与数值求解}
\label{sec:diff-sampling}
\subsection{祖先采样的输入、步骤与输出}
祖先采样（ancestral sampling）从生成联合分布的根节点开始，依次按条件分布抽样。输入是训练好的网络、噪声日程、反向方差以及所需条件；初始化 $X_T\sim p_T$，然后按递减时间使用式\eqref{eq:diff-eps-mean}求均值，并抽取 $X_{t-1}=\mu_\theta(X_t,t)+\gamma_tZ$。各步新增的标准高斯应独立于此前生成状态所用的随机数，最后按观测模型输出 $X_0$ 或约定的重建值。终止条件是抵达指定的最低噪声或零时刻，而非某次网络损失变小。

该流程对选定的离散生成模型可以是精确的条件抽样，但选定模型本身可能只是数据分布的近似。若每步网络成本为 $C$，执行 $K$ 次网络评价的大致计算量为 $KC$，还需考虑引导是否需要双分支、网络批处理和其他附加计算。步数不总等于网络评价次数（number of function evaluations, NFE）：高阶方法一次宏观步可能求值多次，多步法也可能复用历史结果。公平比较应报告实际求值和延迟。

\subsection{DDIM 如何允许确定性路径}
去噪扩散隐式模型（denoising diffusion implicit model, DDIM）表明，可在保持相关单时刻扰动结构与学习目标联系的同时改变路径耦合\cite{song2021ddim}。对选定的两个时刻 $0\leq s<t$，先从 $x_t$ 估计 $\widehat x_0$ 和 $\widehat\varepsilon$，再采用
\begin{align}
 x_s&=a_s\widehat x_0+\sqrt{b_s^2-\sigma_{t\to s}^2}\,\widehat\varepsilon+\sigma_{t\to s}z,
 \label{eq:diff-ddim}\\
 \sigma_{t\to s}^2&=\eta^2\frac{1-\bar\alpha_s}{1-\bar\alpha_t}
 \left(1-\frac{\bar\alpha_t}{\bar\alpha_s}\right),\qquad 0\leq\eta\leq1,
 \label{eq:diff-ddim-var}
\end{align}
其中 $z\sim\mathcal N(0,I_d)$，且 $a_0=1,b_0=0$。当 $\eta=0$ 时没有新增随机项，路径在给定初始噪声后确定；当使用相邻时刻且 $\eta=1$ 时，上述方差和均值恢复相应的后验方差式 DDPM 更新。

跨越时间点必须使用跨越后的累计系数，不能把原来相邻一步的 $\beta_t$ 直接套在任意跨度上。确定性 DDIM 保留了由起始噪声决定的多样性，并可在合适极限下与某种 ODE 离散化联系起来；但是有限大步长下，它并不因噪声网络最优就自动成为原数据分布的精确采样。条件均值预测仍可能混合多个解释，离散更新也有近似误差。把 $\eta$ 调为零是改变采样过程，不是证明扩散本质上无需随机性。

\subsection{欧拉法、二阶法与半线性结构}
对 ODE $\dot x=b(x,\tau)$，欧拉法用当前位置速度预测下一点；在标准光滑性和稳定性条件下，其单步局部误差为步长平方量级，固定区间全局误差通常为步长一阶。二阶预测校正方法先作预测，再使用新位置速度修正，可提高光滑问题上的精度，但通常需要更多求值。对于反向 SDE，Euler--Maruyama 方法还要加入标准差为 $g\sqrt{h}$ 的高斯增量；确定性 ODE 的阶数不能未经解释直接用于随机路径误差。

扩散 ODE 经常具有已知线性项加神经网络非线性项的半线性结构。若把线性项通过积分因子精确处理，再近似剩余积分，可能比直接对全部向量场使用通用欧拉法更有效。DPM-Solver 属于利用这种结构的专用方法\cite{lu2022}。其理论阶数依赖解和网络输出关于所用时间坐标的正则性，文献标题中的少步结果是具体实验表现，不是对任意分布、任意引导强度和任意网络的统一保证。

步长按物理时间均匀、按噪声标准差均匀、按对数信噪比均匀，会把求值点放在不同区域。噪声很小时，向量场可能变化快，端点公式也可能不稳定；较强引导可能进一步放大曲率或刚性。自适应求解器根据局部误差估计调节步长，但其误差控制只针对当前近似向量场，不会消除网络得分的偏差。使用更精确的求解器可能更准确地求解一个有偏模型，因此样本质量不一定随数值容差单调改善。

\subsection{采样误差的四个来源}
完整生成误差至少应分成终点先验失配、学习误差、数值积分误差以及输出处理误差。裁剪预测像素、动态阈值或把最终均值作为输出，会改变模型指定的映射；它们可以是合理的工程选择，却不是无代价的数学恒等变换。尤其在潜空间，坐标通常不天然限制在像素范围，直接套用像素裁剪可能破坏已学习分布。讨论某种输出修正时，应先说明数据尺度和它所改变的对象。

提高步数主要针对离散误差，增加训练数据主要针对统计估计，改进终点日程针对先验失配，重新校准编码与解码针对表示误差。若故障是文本条件被忽视，单纯增加采样步数未必有用；若网络预测准确而粗步长造成细节损坏，重新训练巨大模型也不是首要措施。把误差来源分开，才能提出有针对性的改进和可以解释的实验，而不是把全部变化都归为“模型能力”。

\section{去噪网络的前置知识与结构原理}
\label{sec:diff-network}
\subsection{神经网络学习的是时间条件向量场}
前面的数学目标要求一个从 $(x,t,c)$ 到向量的函数，其中 $x$ 是当前受扰状态，$t$ 是噪声条件，$c$ 是可选外部条件。对噪声预测，输出与输入具有相同数据维数；对得分预测，输出是每个坐标上的对数密度梯度估计。隐藏层可以改变通道数和空间分辨率，但最终必须恢复所需维度。一个输入既可能在低噪声时表示真实结构，也可能在高噪声时只是偶然噪声，因此网络若不知道噪声尺度，就会把不同条件任务混在一起。

神经网络通常是仿射映射与非线性函数的复合，参数通过链式法则计算梯度。反向传播（backpropagation）是高效计算损失对参数导数的方法，与反向扩散的“反向”含义不同：前者沿计算图传播导数，后者沿噪声时间生成状态。一次常见的扩散训练迭代只对一个随机时间的预测做反向传播，并不要求穿过整条采样链。除非明确采用端到端轨迹目标或蒸馏设计，否则把全部生成步骤纳入训练计算图会引入不必要的成本。

不同时间共享参数，可以让网络复用边缘、形状和语义表示；同时也形成多任务冲突。高噪声任务更依赖全局先验，低噪声任务更依赖局部残差，这种分工是有用的解释，但不应被理解成所有样本严格遵循固定的“先语义后纹理”阶段。噪声日程、网络感受野、条件强度和数据模态都会影响实际信息恢复顺序。可靠论述应将这种经验性层次与由系数决定的信噪比区分开。

\subsection{卷积、多尺度结构与跳跃连接}
卷积（convolution）通过共享局部滤波器处理图像，利用平移附近的统计相似性。多层卷积扩大感受野，下采样让较少位置汇总较大空间范围，随后上采样恢复分辨率。U 形网络（U-Net）把下采样路径与上采样路径连接起来，使低层空间细节与高层上下文同时参与预测。跨尺度跳跃连接提供定位信息，残差连接则让一个模块学习对已有表示的修正；二者都叫“跳跃”，但连接的对象和目的不同。

加噪图像并不是每个位置独立可判定的。一个局部斑点是笔画、纹理还是噪声，往往取决于全局布局；而最终精细修正又需要保留位置信息。多尺度网络因此同时需要全局汇聚与局部恢复。高分辨率层计算量较大，低分辨率层适合处理大范围关系；设计时可在不同尺度安排不同宽度、注意力或残差块。这样的选择主要改变表示和计算预算，不改变前向概率模型的基本定义。

归一化层（normalization layer）控制激活尺度，但不同方法使用不同统计量。批归一化依赖批量统计，组归一化在单个样本内部对通道分组，层归一化则按指定特征维度处理。扩散训练中一个批量可以混合不同噪声尺度，因而使用哪些轴统计、是否依赖其他样本，都可能影响表示。不能仅因名称中有“归一化”就认为它们等价，也不能把隐藏层归一化当成数据概率密度的归一化。

\subsection{时间嵌入与条件调制}
时间嵌入（time embedding）把标量时间或噪声尺度映射为多维向量，常见做法是使用多个频率的正弦和余弦，再经可学习变换。不同频率帮助网络区分细微尺度差异；直接使用时间索引、标准差或对数信噪比，实际上选择了不同的坐标。若训练和推理对同一数字赋予不同噪声含义，网络收到的条件会错位，即使采样公式其他部分都正确也可能失败。

条件调制可以通过加性偏置、缩放和平移实现。若特征为 $h$，条件嵌入产生 $r(c,t)$ 与 $b(c,t)$，可令 $h'=r(c,t)\odot\operatorname{Norm}(h)+b(c,t)$，其中 $\odot$ 表示逐分量乘法。维度必须与被调制的特征轴相容。条件通过多个层次进入网络，能影响不同尺度的处理，而不是仅在输入端拼接后期待网络自行保存。调制参数是函数输出，不是为每个采样步骤单独训练一套网络权重。

\subsection{注意力、文本条件与扩散变换器}
注意力（attention）通过查询、键和值建立内容依赖的信息交换。设图像查询矩阵 $Q\in\mathbb R^{n\times h}$，条件键 $K\in\mathbb R^{m\times h}$，条件值 $V\in\mathbb R^{m\times h_v}$，则
\[
 \operatorname{Attn}(Q,K,V)=\operatorname{softmax}\!\left(\frac{QK^\mathsf T}{\sqrt h}\right)V,
\]
其中 softmax 按每个查询所对应的条件位置归一化。输出每行是条件值的加权组合。自注意力的查询、键、值来自同一序列，交叉注意力（cross-attention）则可让图像位置查询文本词元；后者是注入条件的一种架构方式，不等于概率意义上的分类器引导。

对一行分数 $z_1,\ldots,z_m$，softmax 的第 $j$ 个权重为 $\exp(z_j)/\sum_k\exp(z_k)$，因此权重非负且和为一。点积分数除以 $\sqrt h$ 是一种控制维数引起尺度增长的设计；在分量近似独立、方差适当标准化的解释模型中，未缩放点积的方差随 $h$ 增长。真实网络未必严格满足这些独立假设，所以它是有动机的尺度选择，不是对所有注意力表示的概率定理。

扩散变换器（diffusion transformer, DiT）把图像或潜变量切成块，映射为词元，用变换器（Transformer）处理全局关系，再还原噪声或速度输出\cite{peebles2023}。它改变的是去噪网络骨干，扩散损失和采样逻辑仍可沿用。块越小，词元越多，通常保留更细的局部信息但增加注意力成本；全注意力的两两关系数随词元数平方增长。位置编码让模型区分排列，时间和类别可以通过调制进入模块。是否优于卷积结构取决于数据、规模和预算，不能由架构名字直接推出。

\section{训练组织、优化稳定性与诊断思路}
\label{sec:diff-training}
\subsection{把期望目标落实为训练实验}
一个完整训练实验应先确定数据表示、扰动过程、预测目标、时间分布和损失权重。每次从训练集抽取干净样本，从指定时间分布抽取噪声等级，再独立抽取标准高斯，构造 $x_t=a_tx_0+b_t\varepsilon$，交给网络预测相应目标，最后按所定义权重计算平均损失并更新参数。这里是自然语言的数学流程；目标是什么、随机性来自哪里、期望针对哪些变量，都必须在实验记录中明确，而不能只写“训练一个扩散模型”。

如果进行类别或文本条件训练，条件与样本必须来自同一个数据对；如果进行条件丢弃，应明确丢弃概率以及空条件的表示。数据增强也会改变训练分布。水平翻转对普通物体图像可能合理，但对带文字、手性或方向意义的图像未必保留语义；裁剪可能破坏文本与图像的对应。增强的合法性由任务决定，而不是生成模型可以自动忽略所有标签和内容约束。

训练集、验证集和测试集应按数据来源和近重复关系合理划分。若同一图像的小改动跨集合出现，验证指标可能高估泛化。生成模型还可能记忆稀有或重复样本，因此评价不能只看是否能生成逼真的少量图片。对数据清理和采样权重的改变，应记录它们如何改变目标总体：对某些类别过采样可以改善覆盖，但生成分布也会更接近重加权后的数据分布，而非原始频率。

\subsection{学习率、梯度尺度与指数滑动平均}
学习率决定每次沿估计梯度移动多远。小批量梯度噪声、不同时间的目标尺度和高维误差求和都会影响合适的步长。梯度裁剪（gradient clipping）限制异常大更新，有助于应对偶发不稳定，但会改变实际更新方向或幅度；它不能替代修复错误方差、错误权重或数据异常。若持续大量梯度被裁剪，应检查训练定义而不仅是继续降低阈值。

指数滑动平均（exponential moving average, EMA）用 $\bar\theta_k=\rho\bar\theta_{k-1}+(1-\rho)\theta_k$ 平滑参数，其中 $0\leq\rho<1$。它可减少优化轨迹的短期波动，常作为另一组用于评价或采样的参数。参数平均的网络输出不等于各网络输出的平均，因为神经网络关于参数通常是非线性的；EMA 也不是贝叶斯后验平均。比较结果时应注明使用即时参数还是平均参数，以及平均窗口的实际时间尺度。

低精度计算节省显存和计算，但接近零的累计方差、极端对数信噪比或很大的损失权重容易发生数值问题。累计乘积可以在对数域分析其稳定性，标准差与方差要分开处理，时间边界要保持与公式一致。这里的关键是数值量的含义，而非某种固定实现技巧。若所有时间都出现异常，应先检查数据尺度；若仅端点异常，则优先检查系数消去、小分母和目标转换。

\subsection{按噪声尺度拆解损失}
总体损失下降可能掩盖局部失败。可以按对数信噪比区间查看验证回归误差，分别评估高噪声先验判断、中间结构恢复和低噪声修正。如果改变了时间权重，未加权平均误差与加权训练目标可能朝不同方向变化，这本身不矛盾。更有意义的是比较相同定义的指标，并判断改善发生在哪些输入区域，以及它是否传递到采样后的分布质量。

一个基础预测器可以只输出零噪声、线性高斯条件均值或简单的恒等去噪结果，用于建立量纲正确的参照。对于标准高斯数据，解析最优噪声估计为 $b_tx_t$，不是零；其误差可精确计算。因此在可解小问题上检查学习目标，比只看高维图像是否“像样”更容易发现信号系数、时间索引和归约常数错误。基础算例的作用是验证数学定义，不是用玩具结果证明真实任务效果。

\subsection{从症状回到可能原因}
若输出始终近似纯噪声，可能是时间条件不匹配、噪声符号错误、模型未学习或终点尺度不对；若出现整体过亮或过饱和，可能与数据归一化、引导强度、预测转换或裁剪有关；若样本过于相似，应检查数据覆盖、条件分布、采样随机性和训练记忆。这里列举的是待检验假设，不是凭现象就能唯一诊断的因果结论。每次应固定其他因素，改变一个可以解释的环节。

训练预算应同时报告处理的样本量、优化步数、分辨率、模型规模与时间采样设置。仅比较训练轮数可能不公平，因为数据规模不同；仅比较参数量也忽略不同结构的计算与内存差异。停止训练可以依据独立验证目标、生成评价稳定性以及预算约束，而不能宣称扩散去噪误差必须趋于零。条件方差给出不可消除的总体误差，继续追求零训练误差还可能增加记忆风险。

\section{条件生成与引导机制}
\label{sec:diff-guidance}
\subsection{条件分布是目标的一部分}
无条件模型学习 $p(x_0)$，条件模型学习 $p(x_0\mid c)$，其中 $c$ 可以是类别、文本、低分辨率图像、深度图或其他观测。相同条件通常允许许多正确结果，因此条件生成仍然是分布建模，而不是求一个唯一答案。训练时可令 $\varepsilon_\theta(x_t,t,c)$ 预测噪声，前向扰动机制保持不变，只把条件附加到回归输入。理想最优预测为 $\mathbb E[\varepsilon\mid X_t=x_t,C=c]$，相应得分是条件边缘 $p_t(x_t\mid c)$ 的梯度。

条件输入的表达能力与条件目标的正确性是两回事。文本编码器可以提供丰富表示，但训练图文对的错误和歧义仍会传播到模型。交叉注意力能建立位置与词元的联系，却不保证复杂空间关系、否定或计数都被精确满足。对条件控制的评价，应使用与任务语义相符的检验，而不仅是看总体逼真度。一个生成结果可能很逼真但不符合提示，也可能满足局部条件却缺少多样性。

\subsection{分类器引导的贝叶斯分解}
贝叶斯公式在每个噪声时间给出 $p_t(x\mid c)\propto p_t(x)p_t(c\mid x)$。对 $x$ 求导得
\begin{equation}
 s_t(x\mid c)=s_t(x)+\nabla_x\log p_t(c\mid x).
 \label{eq:diff-classifier-guidance}
\end{equation}
第一项来自无条件得分模型，第二项可由在相应噪声图像上训练的分类器提供。分类器必须理解当前噪声强度；只在干净图片上训练的分类器直接处理高噪声状态，所得梯度未必近似所需条件概率。连续条件可用条件密度或其他可微观测模型，但归一化和建模假设也要明确。

在一个小反向高斯步中，若基础转移近似为 $\mathcal N(x;\mu,\Sigma)$，再乘以条件似然，在线性化 $\log p(c\mid x)$ 后有 $\log p(c\mid x)\approx\log p(c\mid\mu)+g^\mathsf T(x-\mu)$，其中 $g=\nabla_x\log p(c\mid x)|_{x=\mu}$。高斯配方得到均值偏移 $\mu+\Sigma g$，协方差保持 $\Sigma$。这是局部一阶近似，条件似然弯曲明显或步长很大时并不精确。它解释了为什么引导强度应与当前不确定性和步长相联系，而不是随意给图像加固定梯度。

分类器梯度可能追求模型可识别的特征，而不完全符合人的语义判断；条件概率估计失真也可能导致不合理方向。提高梯度系数会放大这种偏差，同时改变采样动力学。即使分类器准确，强化某个条件也会重新分配质量，可能抑制边缘情形和稀有模式。因此引导是一种受控偏置机制，是否值得采用取决于条件符合度、多样性和计算成本的共同要求。

\subsection{无分类器引导与系数约定}
无分类器引导（classifier-free guidance, CFG）通过同一个网络学习有条件与无条件预测\cite{ho2022cfg}。训练时以一定概率随机丢弃条件，使空条件分支近似无条件任务；生成时组合两个输出。本文统一采用
\begin{equation}
 s_{\mathrm{cfg}}=s_u+w(s_c-s_u),\qquad
 \varepsilon_{\mathrm{cfg}}=\varepsilon_u+w(\varepsilon_c-\varepsilon_u),
 \label{eq:diff-cfg}
\end{equation}
其中 $u$ 表示空条件，$c$ 表示指定条件。$w=0$ 为无条件，$w=1$ 为普通条件预测，$w>1$ 为外推增强。另一些论述使用 $s_c+\omega(s_c-s_u)$，两种记号的关系是 $w=1+\omega$；比较“引导系数为一”时必须先检查约定。

由于噪声与得分之间是同一时间下的线性缩放，式\eqref{eq:diff-cfg}的两种组合相容。空条件不一定表示空字符串本身的自然语言含义，而是训练时指定的缺失条件状态；如果训练与推理使用的空条件表示不同，无条件分支的解释会失效。条件丢弃若依赖样本内容，还可能使空条件分支学习一个选择偏置分布，因此常见理论解释默认丢弃机制与数据内容适当独立。

\subsection{幂次密度解释及其边界}
若两个得分精确，在固定时间上可以形式化地写成
\[
 s_{\mathrm{cfg}}(x,t)=\nabla_x\log\left[p_t(x\mid c)^w p_t(x)^{1-w}\right].
\]
只要右侧可归一化，它是一个重新加权密度的得分，说明 $w>1$ 倾向强调条件相对无条件更支持的区域。但这些逐时间重新加权密度通常不构成原前向 SDE 的同一条一致边缘路径。因此不能仅凭这个逐点恒等式，就断言有限步引导采样最终精确服从 $p_0(x\mid c)^wp_0(x)^{1-w}$；还需要动力学一致性和求解条件，实践中通常不具备这种直接保证。

引导尺度过大时，条件差分中的误差被放大，向量场可能离开训练时经常覆盖的区域，造成过饱和、结构畸变或多样性下降。重缩放、限制或按时间调节引导，是对这种现象的工程应对，但每种调整都改变实际采样场。若使用负提示代替无条件分支，参考分布也随之改变，已不能直接解释为训练数据的无条件边缘。报告结果应说明两个分支各自输入的条件、系数约定和是否使用额外修正。

\subsection{多种控制信号与可识别性}
将文本、布局和图像条件同时输入，可以让模型近似联合条件分布 $p(x\mid c_1,c_2)$；这需要训练分布支持这些条件的对应关系。把独立训练的条件得分简单相加，有时可以构造有用的组合，但其正确性往往依赖条件独立或乘积专家模型等假设。若两个条件冲突，采样器不会自动知道应优先满足哪一个。必须定义目标偏好，或者承认对应条件事件在训练分布中概率很低。

从生成结果看起来满足某个属性，并不能识别模型内部是否真正理解该属性。它可能利用与属性相关的纹理或背景捷径。可通过保持其他因素不变、改变条件中的关系和计数、检查罕见组合等方式评价控制能力。这里的重点是让评价与声称的能力对齐：若主张空间关系理解，就需要检验关系；若主张文字渲染，就需要检验字符内容，不能用总体图像相似度替代。

\section{潜空间扩散与自编码器的作用}
\label{sec:diff-latent}
\subsection{为什么把扩散移到表示空间}
高分辨率图像的像素维度很大，多步去噪成本随空间尺寸增长。潜空间扩散（latent diffusion）先用编码器 $E$ 将图像映射为潜变量 $Z_0=E(X_0)$，在潜空间学习扩散分布，最后通过解码器 $D$ 生成 $D(Z_0)$\cite{rombach2022}。若编码器是随机的，则 $Z_0\sim q_\phi(z\mid X_0)$，训练潜变量分布包含编码器随机性。这里压缩了表示空间，但压缩质量与生成质量之间存在新的权衡。

设原图大小为 $H\times W\times C$，空间下采样因子为 $r$，潜通道数为 $C_z$，则潜变量坐标数约为 $(H/r)(W/r)C_z$。它相对像素坐标数的比例是 $C_z/(r^2C)$，不能只凭下采样因子就声称成本降低为 $r^2$ 倍；网络宽度、注意力和解码成本也会变化。潜空间扩散的优势通常来自把昂贵的多步处理放在较小表示上，而编码与解码只执行少数次。

\subsection{自编码器、VAE 与聚合后验}
普通自编码器通过重建损失学习编码和解码；变分自编码器（variational autoencoder, VAE）还对随机潜变量与先验之间的关系施加变分约束。训练数据经过编码器所形成的总体潜分布称为聚合后验（aggregated posterior），即 $q_Z(z)=\int q_\phi(z\mid x)\,\mathrm dP_{\mathrm{data}}(x)$。它一般不因为每个样本都受到某种 KL 正则就精确等于标准高斯。潜扩散学习的是这一实际潜分布或其特定缩放版本，而不是根据编码器名称假定它已经可直接抽样。

重建损失可以使用像素误差、特征空间误差或其他目标，不同选择保留的信息不同。像素均方误差重视精确数值，感知特征误差可能更重视语义与纹理，额外的对抗目标可以影响局部逼真度。本文不把这些目标视为可互换的质量定义。若编码器已丢失小文字、细线或计量细节，再准确的潜扩散也无法保证逐样本恢复原来丢失的信息；解码器可能补出合理细节，但这与忠实保留不同。

\subsection{潜变量尺度、噪声日程与坐标变换}
若潜变量每维方差很大，却仍使用为单位尺度设计的噪声日程，高噪声端可能残留显著信号；若方差很小，早期加噪就可能过快破坏结构。因而潜变量的固定缩放影响有效信噪比，应与训练和采样保持一致。设使用 $\widetilde Z_0=cZ_0$，生成后必须以 $D(\widetilde Z_0/c)$ 解码。缩放系数错位会改变解码器输入分布，不是靠多运行几个去噪步骤就能自动补偿。

更一般的线性坐标变换会把各向同性噪声变成非各向同性噪声。若 $Y=AZ$ 且 $A$ 可逆，密度变换含雅可比行列式，得分按 $s_Y(y)=A^{-\mathsf T}s_Z(A^{-1}y)$ 转换。对于非线性且降低维数的编码器，通常不存在简单的可逆换元公式。这说明潜空间得分不是像素空间得分的直接同义词，潜空间上的标准高斯扰动也不等于像素上的标准高斯扰动。

\subsection{两阶段误差和端到端分布}
假设确定性解码器 $D$ 是 $L$-Lipschitz，即 $\|D(z)-D(z')\|\leq L\|z-z'\|$，且相关二阶矩有限。对任何潜分布 $P,Q$，用同一耦合解码可得 $W_2(D_\#P,D_\#Q)\leq L W_2(P,Q)$，其中 $D_\#$ 表示推前分布。再用三角不等式，最终生成分布与真实数据之间的距离可由重建分布误差加上潜分布学习误差经解码放大的项控制。这是一个有条件的分析框架，实际解码器的有效常数可能很大，不能未经估计就作数值保证。

该分解说明应分别评价自编码器重建和潜生成。若重建阶段已存在明显结构缺陷，首先改进表示；若重建良好但随机生成差，则检查潜分布学习与采样；若仅特定条件失败，则检查条件数据和注入机制。最终分布还受到解码器随机性、输出裁剪和颜色变换影响。只在潜空间计算训练损失，不能完整代表用户看到的图像质量。

\subsection{潜空间建模的适用范围}
潜扩散适合能容忍适量表示压缩、同时需要高分辨率生成的任务。对于要求精确数值、微小结构或严格测量一致性的应用，表示损失可能成为主要限制；此时应选择保真编码器、增加相关监督，或直接在观测空间处理。模型复杂度降低与信息保留之间不存在免费的普遍改进。自编码器、扩散先验和解码器组成一个整体系统，任何一个模块的数据偏差都可能在最终结果中显现。

\section{逆问题、图像恢复与后验抽样}
\label{sec:diff-inverse}
\subsection{先定义观测过程，再定义恢复任务}
逆问题（inverse problem）从不完整或受扰观测推断隐藏对象。以 $Y=AX_0+\xi$ 为例，$A\in\mathbb R^{m\times d}$ 是已知观测算子，$\xi\sim\mathcal N(0,\sigma_y^2I_m)$ 为独立观测噪声，$\sigma_y>0$。图像修复的 $A$ 可以选择已知像素，降采样的 $A$ 可以表示模糊后抽取，某些成像问题则使用不同线性算子。若观测机制是非线性的，应写成 $Y=\mathcal A(X_0)+\xi$，不能继续无条件使用线性公式。

贝叶斯恢复目标为 $p(x_0\mid y)\propto p_{\mathrm{data}}(x_0)p(y\mid x_0)$。先验偏好合理数据，似然约束观测一致性；后验抽样保留多种符合观测的解释。若只需要一个点估计，应明确是后验均值、最大后验还是某个任务损失下的 Bayes 决策。扩散先验可以提供候选图像分布，但不能从信息不足的观测中唯一识别本来不可识别的细节。生成的细节即使逼真，也未必等于真实缺失内容。

\subsection{带噪状态上的似然不是原始似然}
在扩散时间 $t$，条件得分满足 $s_t(x_t\mid y)=s_t(x_t)+\nabla_{x_t}\log p_t(y\mid x_t)$，但其中
\begin{equation}
 p_t(y\mid x_t)=\int p(y\mid x_0)q(x_0\mid x_t)\,\mathrm dx_0.
 \label{eq:diff-noisy-likelihood}
\end{equation}
观测由干净数据产生，因此必须对给定受扰状态时的干净数据后验边缘化。直接把 $x_t$ 放入 $p(y\mid x_0)$，相当于假设观测来自受扰图像，改变了问题。即使观测模型很简单，这个积分仍可能困难，因为真实数据先验高度复杂。

一种常见近似是以去噪估计 $D_\theta(x_t,t)$ 替代未知 $x_0$，构造 $\log p(y\mid D_\theta(x_t,t))$ 的梯度。对高斯观测，链式法则产生
\[
 \nabla_{x_t}\log p(y\mid D_\theta)=
 \frac1{\sigma_y^2}J_{D_\theta}(x_t,t)^\mathsf T A^\mathsf T[y-AD_\theta(x_t,t)].
\]
它忽略了条件分布的完整不确定性，即以非线性函数作用于均值代替先积分再取对数，通常不精确。尤其高噪声阶段，$q(x_0\mid x_t)$ 很宽，代入均值的似然可能严重高估确定性。有关方法的性能必须结合这种近似边界解释。

\subsection{线性高斯问题提供精确参照}
若先验为 $X_0\sim\mathcal N(m,C)$，$C$ 正定，则后验仍为高斯：
\begin{align}
 C_{\mathrm{post}}&=(C^{-1}+A^\mathsf T A/\sigma_y^2)^{-1},\nonumber\\
 m_{\mathrm{post}}&=C_{\mathrm{post}}(C^{-1}m+A^\mathsf T y/\sigma_y^2).
 \label{eq:diff-linear-inverse}
\end{align}
这可由第\ref{sec:diff-gaussian}节的精度相加直接推出。即便 $A$ 不满秩，正定先验也使该矩阵可逆；但后验在观测未约束方向仍保留不确定性。当观测噪声趋于零时，需要通过极限或受约束高斯分析处理，不应把零直接代入分母。

在这个可解问题上，正确算法不仅要给出接近后验均值的结果，还应再现后验协方差。一个方法若所有样本都极接近 $m_{\mathrm{post}}$，重建误差可能不错，却不代表后验抽样正确。可以比较不同方向的样本方差、观测残差分布和相关性，检验算法是否过度收缩。这种验证区分了“有用的恢复器”与“校准的后验采样器”，二者目标相关但并不相同。

\subsection{修复、编辑与约束的一致性}
图像修复可把已知区域作为条件，并在适当噪声层级上施加观测约束。若每一步都把已知区域直接替换成完全干净像素，而未知区域仍处于高噪声状态，联合状态可能不符合模型训练时的噪声分布。对已知区域加入匹配时间的扰动可以缓解尺度不一致，但仅做局部替换也不自动保证整个联合后验精确，因为跨区域依赖仍需要正确处理。任何“保持已知区域”的方案都应说明约束在干净空间还是受扰空间施加。

图像编辑通常从输入图像先加噪到某个中间时间，再在新条件下反向生成。较小扰动一般保留更多输入信息，较大扰动给模型更大改变空间，但这只是可解释的权衡，不保证每次语义变化单调。编辑与从先验生成不同：起点受到原图条件限制。若要比较编辑算法，应同时衡量目标条件的满足、无关内容的保持和多样性，而不是只评价最终图像的逼真度。

\section{流匹配、随机插值与整流流}
\label{sec:diff-flow}
\subsection{从去噪问题转向速度学习}
流匹配（flow matching）直接学习把一个简单分布输运到目标分布的时间依赖速度场\cite{lipman2023}。其基本问题是寻找 $u_\theta(x,r)$，使 ODE $\mathrm dX_r/\mathrm dr=u_\theta(X_r,r)$ 从噪声分布出发到达数据分布。这里另用 $r\in[0,1]$，并约定 $r=0$ 为噪声、$r=1$ 为数据，方向与前面的前向扩散时间相反。速度与得分相关但不是同一个对象；指定同样的边缘路径，也可能存在不同的输运速度场。

取随机噪声 $Z\sim P_Z$ 与随机数据 $X\sim P_{\mathrm{data}}$，先选择二者的联合耦合。最简单选择是独立抽样，然后构造直线插值 $X_r=(1-r)Z+rX$，每对端点的条件速度为 $U=X-Z$。这条随机插值的边缘密度记作 $p_r$，可用回归目标
\begin{equation}
 \mathcal L_{\mathrm{FM}}=\mathbb E_{r,Z,X}\|u_\theta((1-r)Z+rX,r)-(X-Z)\|^2
 \label{eq:diff-flow-loss}
\end{equation}
学习速度。时间分布与损失权重仍需指定，端点系数与 DDPM 的方差保持插值也不相同。

\subsection{条件速度为何产生正确边缘输运}
在二阶可积的条件下，命题\ref{prop:diff-regression}给出最优速度 $u^*(x,r)=\mathbb E[X-Z\mid X_r=x]$。为说明它为何有意义，取紧支撑光滑测试函数 $\phi$，并假定可交换微分与期望，则
\[
 \frac{\mathrm d}{\mathrm dr}\mathbb E[\phi(X_r)]
 =\mathbb E[\nabla\phi(X_r)\cdot(X-Z)]
 =\mathbb E[\nabla\phi(X_r)\cdot u^*(X_r,r)].
\]
分部积分得到 $\partial_r p_r=-\nabla\cdot(p_ru^*)$ 的弱形式。在 ODE 与连续性方程解足够良好、初始分布匹配等条件下，沿 $u^*$ 的确定性流便可再现相同边缘分布。这个证明解释了为何无需在训练时求解 ODE，就可以用随机插值提供速度监督。

每一对端点的监督直线都很直，不代表学习到的边缘速度轨迹也必然直。不同端点配对的直线可能交叉；在同一个 $(x,r)$ 上，回归器平均来自不同配对的速度，最终流线可能弯曲。若假设速度场满足局部唯一性条件，真实 ODE 轨迹在相同时刻不会任意交叉，而训练用条件路径却允许交叉。这一区别是理解流匹配的关键：网络学习群体的输运速度，不是在运行时记住某对训练端点并照着其直线走。

\subsection{一般高斯路径与得分的联系}
更一般地设 $X_r=a(r)X+b(r)Z$，$a,b$ 可微且在内部非零，$Z$ 为独立标准高斯。条件速度为 $\dot aX+\dot bZ$，最优边缘速度为
\[
 u^*(x,r)=\dot a\,\mathbb E[X\mid X_r=x]+\dot b\,\mathbb E[Z\mid X_r=x].
\]
由式\eqref{eq:diff-score-identity}，$\mathbb E[Z\mid X_r=x]=-b\nabla\log p_r(x)$，以及 $\mathbb E[X\mid X_r=x]=(x+b^2\nabla\log p_r(x))/a$，得到
\begin{equation}
 u^*(x,r)=\frac{\dot a}{a}x+
 \left(\frac{\dot a}{a}b^2-\dot b\,b\right)\nabla\log p_r(x).
 \label{eq:diff-flow-score}
\end{equation}
这说明在指定高斯路径上，去噪、得分和速度可以相互转换，但转换依赖路径及时间坐标，不能只按网络输出名称机械替换。

某条插值路径能用速度回归学习，不等于它就是某个已指定前向马尔可夫扩散的样本路径。要建立与 SDE 的关系，还需检查其边缘方差、漂移和扩散系数是否相容。概率流 ODE 是从特定 SDE 边缘推导而来，流匹配则可以先指定更广泛的条件路径再学习速度。两者共享连续性方程语言，研究侧重点却有所不同。

\subsection{整流与端点耦合的重要性}
整流流（rectified flow）关注学习较易积分的输运，并可以利用已学习流诱导的端点配对重新构造训练问题\cite{liu2023rectified}。独立噪声与数据的配对可能绕行较多；如果端点间关系更协调，平均后的速度场可能具有更简单的轨迹。重新配对不是简单地把所有弯曲轨迹一条条压成直线，而是改变用于监督的耦合，再学习一个新的边缘速度场。

最优输运（optimal transport）进一步研究在指定代价下哪种耦合最合适。例如平方欧氏代价偏好平均移动距离较短的配对，但一般数据分布上求解准确耦合本身可能昂贵，有限批量近似也会产生偏差。“直线路径”“整流流”和“最优输运”因此不是三个完全等价的名字。较直的轨迹可能降低数值积分成本，但少步质量还取决于速度估计、端点正则性、条件引导及模型容量，不能由路径设计单独保证。

\section{少步生成、渐进蒸馏与一致性模型}
\label{sec:diff-distill}
\subsection{减少步数究竟改变什么}
减少生成成本有三种不同路线：改进求解器，使同一向量场用较少评价达到可接受误差；改变训练路径，使需要求解的动力学更简单；训练一个新模型，直接模仿原模型的较长时间演化。第三类属于知识蒸馏（knowledge distillation）。它把训练计算用于减少推理计算，因而比较成本时不能只看部署步数，还应计入教师生成监督、学生训练及所需数据。三条路线可以组合，但不能用一个方法的理论结论替代另一个方法的验证。

蒸馏的教师提供的是模型行为，而不是真实数据分布的无误答案。若教师遗漏某些模式或对某条件存在偏差，学生模仿得再准确也可能保留这些缺陷；学生容量不足、训练耦合不当或目标不合适还会引入额外误差。直接用真实数据补充监督可以缓解某些问题，但目标随之改变。应明确学生学的是教师终点、教师一步映射、速度场，还是某个分布距离，这些监督不能互相混同。

\subsection{渐进蒸馏与大步映射}
渐进蒸馏（progressive distillation）的一种代表思路是用学生的一步近似教师的两步，再逐阶段减少所需步数\cite{salimans2022}。设教师从时刻 $t$ 到中间时刻 $s$ 的映射为 $F_{t\to s}$，从 $s$ 到 $r$ 的映射为 $F_{s\to r}$，学生拟合其复合 $F_{s\to r}\circ F_{t\to s}$。这样训练目标直接覆盖较大跨度，区别于把只见过小步局部误差的网络临时强行用于大步采样。

若教师是确定性映射，则同一起点对应明确终点，适合直接回归；若教师包含随机转移，直接对随机终点做平方回归只会学习条件均值，可能丢失条件方差。此时需要明确噪声耦合、分布目标或其他机制。一步学生依然可以通过随机起点生成多样样本，但它是否保持教师在给定起点后的路径随机性，是另一件事。蒸馏设计必须先确定要保留哪种分布关系。

\subsection{一致性函数与边界条件}
一致性模型（consistency model）学习一个把同一条概率流轨迹上的不同状态映射到同一低噪声端点的函数\cite{song2023consistency}。设轨迹在时刻 $\tau$ 的状态为 $X_\tau$，最低噪声时刻为 $\delta>0$，理想函数 $F$ 满足 $F(X_\tau,\tau)=X_\delta$，并具有边界条件 $F(x,\delta)=x$。不同时间上输出一致是路径约束，边界条件则把常值解排除，使函数必须与真实端点相联系。

训练可以利用教师向量场得到相邻轨迹点，也可在适当构造下从数据和噪声直接建立一致性约束。用目标网络或停止梯度分支是训练组织手段，不等于已经证明两个预测对应同一真实轨迹。相邻状态如何耦合、使用何种距离、是否满足边界、时间间隔如何选择，都会影响目标含义。若只最小化任意两个输出的差，而没有恰当边界与数据关系，常数函数就可能成为无用解。

一致性输出在理想条件下表示一个确定性流端点，通常不是 $\mathbb E[X_0\mid X_\tau=x]$。后者对前向随机观测下的可能数据求平均，前者沿指定确定性轨迹选取端点，二者依据的耦合不同。把一致性模型称为“直接去噪”有直观帮助，但在数学上必须保留这种差别。一次或少数次函数评价可以构成生成器，实际质量仍受函数近似、端点平滑和训练覆盖限制。

\subsection{质量、覆盖和速度之间的检验}
少步生成特别需要检查条件强度、样本多样性和异常输入。为固定一套提示词得到的学生，未必在不同引导系数或分辨率下保持教师表现。若学生损失只约束某种教师轨迹，它还可能对更换采样日程敏感。评价应固定计算预算并报告教师与学生的实际延迟、生成覆盖、条件指标及失败案例，而不能只报告网络调用次数减少了多少。

一个合理实验可以同时比较原教师的粗步采样、改进数值求解器、相同参数规模的蒸馏学生以及更小网络的普通训练。这样能分清收益来自求解、监督还是架构变化。若学生需要大量教师样本，训练成本应按预期生成次数摊销；在只生成少量样本的场景中，更贵的训练未必划算。算法选择应对应真实使用频率，而不是把最低推理步数视为唯一目标。

\section{离散扩散与其他数据模态}
\label{sec:diff-discrete}
\subsection{从高斯核变为转移矩阵}
对于类别或词元，直接加高斯噪声后再舍入不一定尊重原始离散空间。离散扩散（discrete diffusion）可以用有限状态马尔可夫核扰动数据\cite{austin2021}。设单个变量取 $K$ 个状态，采用行向量概率约定，转移矩阵 $Q_t(i,j)=\Pr(X_t=j\mid X_{t-1}=i)$ 每行非负且和为一。累计转移为 $\overline Q_t=Q_1\cdots Q_t$，给定 $X_0=k$，时刻 $t$ 取 $j$ 的概率是 $\overline Q_t(k,j)$。矩阵相乘的顺序与行列约定有关，必须先约定再推导。

当分母为正时，贝叶斯公式给出
\begin{equation}
 \Pr(X_{t-1}=i\mid X_t=j,X_0=k)
 =\frac{\overline Q_{t-1}(k,i)Q_t(i,j)}{\overline Q_t(k,j)}.
 \label{eq:diff-discrete-posterior}
\end{equation}
它与连续高斯后验承担相同角色：已知原始样本时可精确计算监督，生成时仍需要对未知原始状态进行统计推断。分母为零表示该条件事件在模型下不发生，不能据此定义任意普通条件概率。离散模型的损失常涉及类别概率的交叉熵或 KL，不能直接把连续噪声均方误差全部搬过来。

\subsection{均匀替换与吸收掩码}
一种扰动是以一定概率保留当前状态，否则替换为随机类别；另一种是逐步把状态变成特殊掩码，并使掩码在前向过程中保持为掩码，形成吸收状态（absorbing state）。反向模型学习在不同掩码程度下恢复词元或类别。使用掩码的好处是噪声语义明确，但转移支持与时间调度不同于高斯平滑，不具有同样的普通欧氏得分。数学对象仍是概率核与后验，而不是对整数编号求导。

对序列各位置独立扰动，不意味着生成时各位置边缘独立。网络可以读取整个受扰序列，联合利用上下文；每一步若把输出条件分解为逐位置类别分布，则仍引入某种条件独立近似。多轮更新能传递依赖，却不自动证明恢复了任意真实联合转移。并行预测减少序列顺序依赖，迭代轮数又增加额外计算；与自回归方法比较时，应同时考虑序列长度、每轮网络成本、缓存能力和任务约束。

\subsection{连续表示中的文本与类别}
也可以先把离散对象映射到连续嵌入，再在连续空间扩散，最后用解码规则恢复类别。这样便于复用高斯机制，但连续空间的距离未必等同于语义关系，最终舍入或解码会改变分布。若许多连续点映射到同一个类别，生成类别概率取决于整个解码区域的质量，而不是某个嵌入中心的密度。评价必须落到最终离散对象上，不能只看嵌入均方误差。

文本生成还涉及变长、语法、逻辑一致性和外部事实等问题。一个模型在去噪中能恢复常见词，并不自动说明其推理能力与另一类生成架构相同。扩散描述的是生成组织方式，任务能力还取决于数据、容量、表示、监督和评价。比较不同生成范式，应提出具体任务和计算约束，避免从概率路径的形式直接推断通用智能水平。

\subsection{音频、视频与几何数据}
音频可以在波形、时频表示或压缩潜空间上建模。波形维度很高且对相位和局部连续性敏感，时频表示又需要处理幅度、相位及重建约束；生成表示正确不保证还原波形同样正确。视频增加了时间轴，网络需要维护跨帧对象、运动和相机关系。逐帧图像都逼真仍可能产生闪烁，因此评价要同时观察单帧质量和时序一致性，而不是简单平均图像指标。

三维点云、分子和姿态包含平移、旋转、排列等对称性。若整体旋转输入，理想坐标预测应相应旋转，这称为等变性（equivariance）；某些整体性质则应保持不变。网络和噪声过程若尊重这些结构，可以减少不必要的学习负担。但欧氏坐标噪声可能破坏键长或约束流形，最终需要恰当表示、投影或几何扩散。任意投影也会改变概率规律，应明确它是模型的一部分还是近似后处理。

无论处理哪种模态，扩散框架都需要四项具体定义：数据空间及其参考测度、已知扰动核、可学习反向对象、与最终任务相符的评价方式。高斯图像扩散提供的是一套可迁移思想，而不是在所有空间中无需调整的万能公式。先明确对象再选择工具，能避免把词元、像素、旋转矩阵和物理测量值当成完全相同的实数数组。

\section{生成模型评价、比较与可靠结论}
\label{sec:diff-evaluation}
\subsection{质量、多样性、条件符合度和概率校准}
生成质量不是一个单一量。逼真度关注样本是否符合真实数据的局部与整体结构，多样性关注是否覆盖不同有效模式，条件符合度关注输入要求是否被满足，概率校准关注出现频率与目标概率是否一致。生成一张非常漂亮的图，不能说明分布覆盖良好；覆盖许多外形，也不一定说明其比例正确。对于逆问题，还需要观测一致性与不确定性校准。评价指标应与研究主张逐一对应。

人类评价可以发现指标忽略的细节，但也受到提示词选择、展示顺序、样本筛选和评审者差异影响。应避免只展示挑选后的最佳样本，同时报告随机样本或明确的采样规则。成对比较需要随机化顺序并控制条件，重复评价有助于估计不确定性。自动指标方便大规模比较，却不应被当作对语义、真实性和应用可用性的全部判断。

\subsection{特征分布距离的原理与限制}
Fr\'{e}chet Inception Distance（FID）先将图像映射到固定特征空间，再用高斯均值和协方差近似真实与生成特征分布\cite{heusel2017}。设二者统计量为 $(m_r,C_r)$ 与 $(m_g,C_g)$，其形式为
\begin{equation}
 \operatorname{FID}=\|m_r-m_g\|^2+
 \operatorname{tr}\!\left[C_r+C_g-2(C_r^{1/2}C_gC_r^{1/2})^{1/2}\right].
 \label{eq:diff-fid}
\end{equation}
这里使用对称半正定矩阵平方根，使表达式在协方差半正定时具有清楚含义。它比较特征的前两阶统计量，不是原始图像分布的完整距离。即使总体均值和协方差相同，不同高阶结构也可能使真实视觉分布不同。

有限样本估计的 FID 有统计偏差与方差，还依赖特征网络、图像缩放、裁剪、样本量和真实参考集。比较时必须保持这些设置一致。另一类核距离，如 Kernel Inception Distance（KID），使用特征上的最大均值差异（maximum mean discrepancy, MMD）估计\cite{binkowski2018}。无偏估计可以减少某类有限样本偏差，但并不意味着零方差或对所有语义敏感；任何固定特征都可能遗漏任务关心的信息。

\subsection{似然为什么不等于感知质量}
似然关注模型为真实观测分配多少概率，感知指标关注选定表示或人的判断。高维数据中的大量低层细节可以显著影响对数密度，却不一定影响人眼感受；反过来，强化引导可能提高可见条件符合度，同时改变概率分布。扩散的变分上界还包含后验间隙，与直接计算的模型似然不是完全相同量。报告“似然改善”时应说明是真实可计算似然、数值估计，还是变分界。

模型在测试数据上似然较高，也不自动排除记忆。有限训练样本附近的概率质量可能被过度集中，特别是大量重复数据存在时。可以在像素和语义特征空间检查最近邻，并结合样本频率、生成提示和具体内容分析是否发生复制；但最近邻距离本身也不能完整证明或否定记忆。对于受严格观测约束的任务，逼真结果尤其不能替代测量证据，应保留模型推断与实际观测之间的边界。

\subsection{对照实验如何隔离机制}
若要研究采样器，应固定训练网络、终点噪声和条件，比较相同 NFE 或实际时间下的效果；若研究训练目标，应固定数据、模型规模和训练预算，再分别考察同一采样器及适配采样器的结果；若研究潜表示，则同时报告重建误差和最终生成指标。引导尺度可以显著改变质量与多样性的权衡，应报告相应曲线或一组明确配置，而非为不同方法偷偷选择不同目标点。

随机种子控制便于成对比较，但一个种子的结果仍可能偶然。对核心指标应使用足够样本并估计波动，模型选择与最终测试应分开。一次实验同时改变网络、数据量、噪声日程和采样步数，即使结果更好，也难以说明哪种机制贡献了收益。消融实验（ablation study）的价值在于拆分因果解释，不是机械地删掉模块后列出许多数值。

\subsection{与其他生成范式的关系}
自回归模型（autoregressive model）把联合概率按变量顺序分解，生成依赖先前变量；变分自编码器通过潜变量和近似后验组织学习；生成对抗网络（generative adversarial network, GAN）通过判别式对抗信号训练生成分布；正规化流（normalizing flow）利用可逆变换与密度换元。扩散则通过噪声尺度上的局部推断构造生成过程。它们可以相互结合，例如潜扩散使用自编码器，概率流连接连续正规化流，蒸馏可以引入其他分布匹配目标。

比较这些范式应关注可计算似然、生成并行度、采样成本、训练稳定性、条件控制和表示约束，而不是宣布某一类在所有方面占优。扩散把一次复杂生成拆为多个较局部的任务，代价是多次函数评价；蒸馏可以转移这项成本，却需要额外训练。概率分解、网络架构和训练监督各自影响系统，因此仅凭生成范式无法推出具体模型的全部能力与限制。

\section{综合例题、练习与解答}
\label{sec:diff-exercises}
\subsection{例题一：把一次去噪的所有对象算清楚}
设 $X_0\sim\mathcal N(0,4)$，$\varepsilon\sim\mathcal N(0,1)$ 独立，$X_t=0.6X_0+0.8\varepsilon$。要求计算受扰方差、最佳干净样本预测、最佳噪声预测、得分，以及干净样本预测的最小均方误差。此题的目的不是代入一个孤立公式，而是检查五个对象是否描述同一个联合模型。

\begin{solve}
受扰方差为 $0.6^2\times4+0.8^2=2.08$。联合协方差给出 $\mathbb E[X_0\mid X_t=x]=2.4x/2.08=15x/13$，以及 $\mathbb E[\varepsilon\mid X_t=x]=0.8x/2.08=5x/13$。得分为 $s_t(x)=-x/2.08=-25x/52$，也等于负噪声预测除以 $0.8$。Tweedie 公式得到 $(x+0.64s_t(x))/0.6=15x/13$，相互一致。最小干净数据误差为后验方差 $4\times0.64/2.08=16/13$，不是零。条件均值系数大于一并不矛盾，因为前向把信号收缩为原来的 $0.6$，去噪同时需要补偿尺度。
\end{solve}

\subsection{例题二：双峰分布中的得分与模式}
沿用例\ref{ex:diff-two-point}，取 $a=2,\sigma=1$。要求求出平滑分布的得分，并解释 $y=0$ 附近的行为。由两个高斯的混合，$p(y)=\tfrac12\mathcal N(y;-2,1)+\tfrac12\mathcal N(y;2,1)$。其密度在两个原始点附近形成结构，但真实峰的位置不能直接凭原始均值认定，需检查导数。

\begin{solve}
由去噪恒等式，$s(y)=2\tanh(2y)-y$，所以 $s(0)=0$。其导数在零处为 $s'(0)=4-1=3>0$，表明对数密度在零处局部向两侧上升，零是局部低谷而非高峰。对于很小的正观测，得分为正；很小的负观测，得分为负。虽然条件均值在零点恰为零，邻近点的得分场仍保留双峰结构。若加大噪声到足够程度，中间低谷可以消失，这解释了平滑尺度对多峰几何的影响。
\end{solve}

\subsection{例题三：高斯最优噪声误差与反向均值}
设 $X_0\sim\mathcal N(0,I_d)$，采用标准 VP 日程。证明最佳噪声预测为 $b_tX_t$，最小噪声回归误差为 $d a_t^2$，并核对 DDPM 均值公式。该例同时说明不同噪声时间的原始损失大小为何不直接代表学习能力。

\begin{solve}
由于 $a_t^2+b_t^2=1$，$X_t$ 为标准高斯，且 $\operatorname{Cov}(\varepsilon,X_t)=b_tI_d$，所以条件均值为 $b_tX_t$，条件协方差为 $(1-b_t^2)I_d=a_t^2I_d$。最小平方误差是该协方差的迹，即 $d a_t^2$。代入式\eqref{eq:diff-eps-mean}，均值为 $(x_t-\beta_tx_t)/\sqrt{\alpha_t}=\sqrt{\alpha_t}x_t$。高噪声时 $a_t$ 较小，噪声本身更容易预测，故损失可以较低；这并不意味着高噪声时更容易确定干净数据。
\end{solve}

\subsection{例题四：参数化转换的损失权重}
设 $a=3/5,b=4/5$，网络使用速度输出 $v_\theta$，且某个样本上速度预测误差为向量 $e$。要求计算由同一输出变换得到的干净数据误差与噪声误差，并说明能否把三种无权损失视为数值相同。

\begin{solve}
由 $X_0=aX_t-bV$ 和 $\varepsilon=bX_t+aV$，误差分别为速度误差的 $-b$ 倍和 $a$ 倍，平方范数为 $16\|e\|^2/25$ 与 $9\|e\|^2/25$。它们相对于速度误差具有不同尺度。若 $a,b$ 随时间变化，这些比例还随时间变化，因此在不同参数化上使用同样的无权平均，会改变时间任务之间的权衡。只有携带对应权重并保持函数转换一致，才可以讨论目标等价。
\end{solve}

\subsection{例题五：概率流在一维高斯族上的验证}
考虑零漂移 SDE，边缘分布为 $\mathcal N(0,S(\tau))$，其中 $S(\tau)=S_0+v(\tau)>0$、$v'(\tau)=g(\tau)^2$。推导概率流速度并求解从 $x_0$ 出发的确定性轨迹，解释为什么该流可以增加分布方差而不用随机噪声。

\begin{solve}
得分为 $-x/S(\tau)$，概率流速度为 $g(\tau)^2x/[2S(\tau)]$。因为 $S'=g^2$，分离变量得到 $X_\tau=\sqrt{S(\tau)/S_0}\,x_0$。若随机初始状态服从 $\mathcal N(0,S_0)$，确定性缩放后方差正好为 $S(\tau)$。SDE 通过新增随机增量增加方差，ODE 通过扩张随机初始群体达到相同边缘；路径显然不同。反向积分则收缩群体，仍不等于所有点沿普通得分上升最终坍缩到零。
\end{solve}

\subsection{例题六：两状态离散扩散的后验}
设状态为 $A,B$，每步以概率 $0.8$ 保持、以概率 $0.2$ 切换，即 $Q=\begin{pmatrix}0.8&0.2\\0.2&0.8\end{pmatrix}$。已知 $X_0=A,X_2=B$，求 $X_1$ 的条件分布。这里既可能先保持再切换，也可能先切换再保持，不能只按最终状态猜测中间状态。

\begin{solve}
累计转移 $Q^2$ 的非对角元为 $0.8\times0.2+0.2\times0.8=0.32$。$X_1=A$ 对应路径概率 $0.16$，$X_1=B$ 也为 $0.16$，归一化后各为 $1/2$。这个例子展示了已知起点和终点仍可能无法确定中间历史。若生成时起点也未知，还需要按 $q(X_0\mid X_2)$ 进一步混合，正如连续高斯扩散中的式\eqref{eq:diff-reverse-mixture}。
\end{solve}

\subsection{练习：检查假设而非只记忆公式}
第一题，设无条件得分 $s_u(x)=-x$，某条件得分 $s_c(x)=-(x-m)$，二者对应方差同为一的高斯。按本文 CFG 约定求组合场对应的固定时间高斯。解答是 $s_{\mathrm{cfg}}(x)=-x+wm$，对应均值 $wm$、方差一的高斯；它说明引导可把均值外推到条件分布中心之外，但这仍只是单时刻计算，不能据此跳过跨时间路径一致性。

第二题，证明在正时间权重下，无限容量去噪回归的逐时间最优函数不改变，并说明有限共享网络为何可以改变答案。证明使用式\eqref{eq:diff-projection}：每个时间的额外误差都是正权重乘以预测与条件均值之差的平方，最小值由该差为零取得。有限共享网络无法同时独立调整所有时间，目标变为在受限函数族中折中，权重决定哪一部分误差更昂贵。

第三题，解释为什么“加入的噪声标签已知”不意味着噪声预测的 Bayes 误差为零。应指出标签是训练者通过一次联合抽样知道的，而网络输入没有包含全部抽样变量；不同潜在组合可产生相同输入。若 $X_0$ 与噪声均为非退化高斯，条件噪声协方差通常严格正，例题三给出了精确值。这是信息限制而非优化失败。

第四题，设两个生成系统训练损失相同，却使用不同的反向方差和不同步长，能否推出输出分布相同？不能。训练损失只是一个期望标量，不唯一决定预测函数；即使预测函数相同，反向转移方差和数值离散化也会改变路径与终点分布。要建立分布等价，需要明确模型、初始分布、时间路径和求解条件，不能用一个损失值代替这些定义。

\subsection{一个可独立完成的研究性学习方案}
可以先选一维高斯、双高斯混合或二维圆环作为已知数据分布，分别写出可计算的扰动密度与得分。先使用解析得分比较随机反演和概率流，以隔离数值误差；再用近似函数代替解析得分，观察学习误差如何改变结果；最后加入条件或线性观测，比较真正后验与近似引导。每一步都应检查样本均值、方差、模式权重和终点分布，而不仅是几条轨迹是否好看。

这一学习方案的价值在于控制未知因素：解析模型中可以知道正确答案，高维真实图像中通常做不到。只有在简单问题上确认概率对象、时间方向和参数化一致，才适合把同样思路扩展到复杂网络。本文中的所有数值例题都是推导示例，并不构成训练实验结果；读者据此开展实验时，应重新记录配置、误差和实际观测，避免把理论可解性误写成实践性能保证。

\begin{thebibliography}{99}
\bibitem{sohl2015} Jascha Sohl-Dickstein, Eric A. Weiss, Niru Maheswaranathan, Surya Ganguli. Deep Unsupervised Learning using Nonequilibrium Thermodynamics. 2015. \url{https://arxiv.org/abs/1503.03585}.
\bibitem{hyvarinen2005} Aapo Hyv\"arinen. Estimation of Non-Normalized Statistical Models by Score Matching. Journal of Machine Learning Research, 6:695--709, 2005. \url{https://jmlr.org/papers/v6/hyvarinen05a.html}.
\bibitem{song2019} Yang Song, Stefano Ermon. Generative Modeling by Estimating Gradients of the Data Distribution. 2019. \url{https://arxiv.org/abs/1907.05600}.
\bibitem{ho2020} Jonathan Ho, Ajay Jain, Pieter Abbeel. Denoising Diffusion Probabilistic Models. 2020. \url{https://arxiv.org/abs/2006.11239}.
\bibitem{song2021sde} Yang Song, Jascha Sohl-Dickstein, Diederik P. Kingma, et al. Score-Based Generative Modeling through Stochastic Differential Equations. ICLR, 2021. \url{https://arxiv.org/abs/2011.13456}.
\bibitem{song2021ddim} Jiaming Song, Chenlin Meng, Stefano Ermon. Denoising Diffusion Implicit Models. ICLR, 2021. \url{https://arxiv.org/abs/2010.02502}.
\bibitem{nichol2021} Alexander Quinn Nichol, Prafulla Dhariwal. Improved Denoising Diffusion Probabilistic Models. 2021. \url{https://arxiv.org/abs/2102.09672}.
\bibitem{karras2022} Tero Karras, Miika Aittala, Timo Aila, Samuli Laine. Elucidating the Design Space of Diffusion-Based Generative Models. 2022. \url{https://arxiv.org/abs/2206.00364}.
\bibitem{lu2022} Cheng Lu, Yuhao Zhou, Fan Bao, et al. DPM-Solver: A Fast ODE Solver for Diffusion Probabilistic Model Sampling in Around 10 Steps. 2022. \url{https://arxiv.org/abs/2206.00927}.
\bibitem{ho2022cfg} Jonathan Ho, Tim Salimans. Classifier-Free Diffusion Guidance. 2022. \url{https://arxiv.org/abs/2207.12598}.
\bibitem{rombach2022} Robin Rombach, Andreas Blattmann, Dominik Lorenz, Patrick Esser, Bj\"orn Ommer. High-Resolution Image Synthesis with Latent Diffusion Models. CVPR, 2022. \url{https://arxiv.org/abs/2112.10752}.
\bibitem{peebles2023} William Peebles, Saining Xie. Scalable Diffusion Models with Transformers. ICCV, 2023. \url{https://arxiv.org/abs/2212.09748}.
\bibitem{lipman2023} Yaron Lipman, Ricky T. Q. Chen, Heli Ben-Hamu, Maximilian Nickel, Matt Le. Flow Matching for Generative Modeling. ICLR, 2023. \url{https://arxiv.org/abs/2210.02747}.
\bibitem{liu2023rectified} Xingchao Liu, Chengyue Gong, Qiang Liu. Flow Straight and Fast: Learning to Generate and Transfer Data with Rectified Flow. ICLR, 2023. \url{https://arxiv.org/abs/2209.03003}.
\bibitem{salimans2022} Tim Salimans, Jonathan Ho. Progressive Distillation for Fast Sampling of Diffusion Models. ICLR, 2022. \url{https://arxiv.org/abs/2202.00512}.
\bibitem{song2023consistency} Yang Song, Prafulla Dhariwal, Mark Chen, Ilya Sutskever. Consistency Models. ICML, 2023. \url{https://arxiv.org/abs/2303.01469}.
\bibitem{austin2021} Jacob Austin, Daniel D. Johnson, Jonathan Ho, Daniel Tarlow, Rianne van den Berg. Structured Denoising Diffusion Models in Discrete State-Spaces. 2021. \url{https://arxiv.org/abs/2107.03006}.
\bibitem{heusel2017} Martin Heusel, Hubert Ramsauer, Thomas Unterthiner, Bernhard Nessler, Sepp Hochreiter. GANs Trained by a Two Time-Scale Update Rule Converge to a Local Nash Equilibrium. 2017. \url{https://arxiv.org/abs/1706.08500}.
\bibitem{binkowski2018} Miko{\l}aj Bi\'nkowski, Danica J. Sutherland, Michael Arbel, Arthur Gretton. Demystifying MMD GANs. ICLR, 2018. \url{https://arxiv.org/abs/1801.01401}.
\end{thebibliography}

\end{document}
